% =============================================================================
% DISPENSA — Funzionamento dei sistemi biologici
% Corso strutturato in 4 capitoli, ciascuno composto da piu' lezioni.
% Ogni lezione e' un file \section dentro la cartella del capitolo,
% incluso qui sotto con \input nel capitolo corrispondente.
% Stile: template-schematico.tex (definizioni con i due punti, liste annidate).
% =============================================================================
\documentclass[11pt, oneside, openany]{book}

% --- Lingua e codifica ---
\usepackage[italian]{babel}
\usepackage[T1]{fontenc}
\usepackage[utf8]{inputenc}
\usepackage{lmodern}

% --- Margini e layout (ottimizzato per lettura a schermo, non per la stampa) ---
\usepackage[
  margin=12mm, headsep=3mm, footskip=6mm,
  includeheadfoot
]{geometry}
\setlength{\headheight}{15pt}
\raggedbottom

% --- Spaziatura e tipografia ---
\usepackage{setspace}
\setstretch{1.15}
\setlength{\parindent}{0pt}
\setlength{\parskip}{4pt}
\usepackage[protrusion=true, expansion=true, babel=true]{microtype}

% --- Titoli capitolo (deve stare prima di hyperref) ---
\usepackage{titlesec}
\titleformat{\chapter}[block]
  {\normalfont\Huge\bfseries}
  {\thechapter\quad}{0pt}{}
\titlespacing*{\chapter}{0pt}{-60pt}{6pt}

% --- Link e segnalibri PDF ---
\usepackage[hidelinks]{hyperref}
\usepackage{bookmark}

% --- Intestazioni e piè di pagina ---
\usepackage{fancyhdr}
\pagestyle{fancy}
\fancyhf{}
\fancyhead[L]{\leftmark}
\fancyhead[R]{\rightmark}
\fancyfoot[C]{\thepage}
\renewcommand{\headrulewidth}{0.4pt}
\renewcommand{\footrulewidth}{0pt}
\fancypagestyle{plain}{%
  \fancyhf{}%
  \fancyfoot[C]{\thepage}%
  \renewcommand{\headrulewidth}{0pt}%
}
\renewcommand{\chaptermark}[1]{\markboth{Cap.~\thechapter:\ #1}{}}
\renewcommand{\sectionmark}[1]{\markright{\thesection.\ #1}}

% --- Matematica e grafica ---
\usepackage{amsmath}
\usepackage{amssymb}
\usepackage{graphicx}

% --- Tabelle ---
\usepackage{booktabs}
\usepackage{array}
\usepackage{tabularx}
\usepackage{multirow}

% --- Liste annidate (strumento principale dello stile schematico) ---
\usepackage{enumitem}
\setlist[itemize]{noitemsep, topsep=2pt}
\setlist[enumerate]{noitemsep, topsep=2pt}

% --- Due colonne per liste lunghe parallele ---
\usepackage{multicol}

% =============================================================================
% COMANDI DI STILE SCHEMATICO
% =============================================================================
% Sotto-titolo in grassetto, non numerato, su riga propria (es. "Struttura", "DNA").
\newcommand{\rubrica}[1]{\par\addvspace{4pt}\noindent\textbf{#1}\par\nobreak}

% Alberi di classificazione: minimali, solo testo + rami (nessun box).
\usepackage{forest}
\forestset{
  classalbero/.style={
    for tree={
      grow'=east, parent anchor=east, child anchor=west,
      font=\footnotesize, anchor=west, align=center, edge={-},
      l sep=8mm, s sep=2mm, inner sep=1pt,
    }
  }
}

% --- Metadati ---
\title{Funzionamento dei sistemi biologici}
\author{Prof. Michela Montorsi}
\date{A.A. 2024/2025}

% =============================================================================
\begin{document}

\frontmatter                 % numerazione romana (i, ii, ...) per titolo e indice

\maketitle
\thispagestyle{empty}
\clearpage

\tableofcontents
\clearpage

\mainmatter                  % la numerazione delle pagine riparte da 1

% =============================================================================
% CAPITOLO 1 — Organismi viventi
% =============================================================================
\chapter{Organismi viventi}
\section{Biologia}

\begin{itemize}
  \item \textbf{evoluzione}: le popolazioni di organismi si sono evolute nel tempo a partire da forme di vita primordiali;
  \item \textbf{trasmissione dell'informazione}: l'evoluzione dipende dalla trasmissione dell'informazione genetica da una generazione all'altra;
  \item \textbf{trasferimento dell'energia}: tutti i processi vitali richiedono un continuo ingresso di energia.
\end{itemize}

% ---------------------------------------------------------------------------
\subsection{I livelli di organizzazione della vita}

Gerarchia dei livelli, dal più semplice al più complesso:
\[
\text{atomo} \rightarrow \text{molecola} \rightarrow \text{macromolecola} \rightarrow \text{procariote} \rightarrow \text{organello} \rightarrow \text{cellula eucariote} \rightarrow \text{tessuto}
\]
\[
\rightarrow \text{organo} \rightarrow \text{apparato} \rightarrow \text{organismo} \rightarrow \text{popolazione} \rightarrow \text{comunità} \rightarrow \text{ecosistema} \rightarrow \text{biosfera}.
\]

\rubrica{I viventi: livelli di organizzazione}
\begin{center}
\begin{forest} classalbero
[I viventi
  [Livello molecolare [Virus][Viroidi][Prioni]]
  [Livello cellulare
    [Procarioti [Archea][Batteri]]
    [Eucarioti [Protozoi][Lieviti][Cellule animali][Cellule vegetali]]
  ]
]
\end{forest}
\end{center}

% ---------------------------------------------------------------------------
\subsection{La classificazione dei viventi}

\textbf{Tre domini} (dall'analisi delle sequenze di rRNA di Carl Woese; tutti derivati da un progenitore comune):
\begin{multicols}{3}
  \begin{itemize}
    \item \textbf{Bacteria} (Batteri);
    \item \textbf{Archaea} (Archei);
    \item \textbf{Eukarya} (Eucarioti).
\end{itemize}
\end{multicols}

\textbf{Sei regni}: Bacteria, Archaea, Protista, Plantae, Animalia, Fungi.

\rubrica{Gerarchia di classificazione linneana}
Dal raggruppamento più ampio al più specifico:
\[
\text{dominio} \rightarrow \text{regno} \rightarrow \text{phylum} \rightarrow \text{classe} \rightarrow \text{ordine} \rightarrow \text{famiglia} \rightarrow \text{genere} \rightarrow \text{specie}.
\]
Esempi:
\begin{itemize}
  \item Eukarya $\rightarrow$ Animalia $\rightarrow$ Chordata $\rightarrow$ Mammalia $\rightarrow$ Primates $\rightarrow$ Pongidae $\rightarrow$ \textit{Pan} $\rightarrow$ \textit{Pan troglodytes};
  \item Eukarya $\rightarrow$ Animalia $\rightarrow$ Artropodi $\rightarrow$ Insetti $\rightarrow$ Blattoidei $\rightarrow$ Blattellidi $\rightarrow$ \textit{Blattella} $\rightarrow$ \textit{Blattella germanica}.
\end{itemize}

\textbf{Specie}: insieme di popolazioni interfertili fra loro e isolate riproduttivamente dalle altre.

% ---------------------------------------------------------------------------
\subsection{Le teorie evoluzionistiche}

\rubrica{Linneo (fine 1700)}
Fondatore della moderna classificazione dei viventi; definiva le specie come ``fisse'', cioè invariabili nel tempo e incapaci di modificarsi. All'inizio del XIX secolo la teoria fu messa in discussione: negli strati rocciosi più antichi mancavano tracce fossili delle specie allora viventi, mentre se ne rinvenivano altre appartenenti a organismi estinti.

\rubrica{Lamarck (1809, \textit{Philosophie Zoologique})}
Prima teoria evoluzionista: ipotizzò la modificabilità delle specie sotto l'influenza delle condizioni ambientali.
\begin{itemize}
  \item principio: è il bisogno di una funzione che crea l'organo; l'uso o il non uso di determinati organi porta con il tempo a un loro potenziamento o a una loro atrofia;
  \item errore di fondo: l'ereditabilità dei caratteri acquisiti.
\end{itemize}

\rubrica{Darwin (1859, \textit{L'origine delle specie per mezzo della selezione naturale})}
Nuova teoria evolutiva basata su tre concetti fondamentali: \textbf{variabilità}, \textbf{selezione naturale}, \textbf{ereditarietà}.
\begin{itemize}
  \item viaggio del \textit{Beagle} (1831--1836): circumnavigazione del globo;
  \item albero della vita: la storia della vita immaginata come un albero; i punti di ramificazione rappresentano l'origine di nuove famiglie, i rami che non raggiungono il margine superiore rappresentano gruppi estinti.
\end{itemize}
Osservazioni e deduzioni:
\begin{enumerate}
  \item la maggior parte degli organismi produce più di uno o due figli; le popolazioni non aumentano indefinitamente di dimensioni; cibo e risorse sono limitati $\rightarrow$ gli individui competono per risorse limitate;
  \item gli individui mostrano variabilità in molte caratteristiche; molte variazioni hanno basi genetiche ereditate $\rightarrow$ le caratteristiche ereditarie permettono ad alcuni individui di sopravvivere più a lungo e riprodursi più di altri;
  \item conclusione: le caratteristiche di una popolazione cambieranno nel corso delle generazioni perché quelle vantaggiose ed ereditabili diventeranno più comuni.
\end{enumerate}

% ---------------------------------------------------------------------------
\subsection{La teoria cellulare}

Tre principi:
\begin{itemize}
  \item le cellule sono le unità fondamentali della vita;
  \item tutti gli organismi sono composti da cellule;
  \item tutte le cellule provengono da cellule preesistenti.
\end{itemize}

\rubrica{Scoperta delle cellule}
\begin{itemize}
  \item \textbf{Robert Hooke} (1665): con un microscopio composto osserva, in una sezione sottile di sughero, un reticolo di ``celle'';
  \item \textbf{Anton van Leeuwenhoek} (1673): con un microscopio a lente singola (ingrandimento $\sim$270 volte, risoluzione $\sim$1,35\,$\mu$m) osserva batteri e altri microrganismi.
\end{itemize}

% ---------------------------------------------------------------------------
\subsection{Ordini di grandezza e microscopia}

\rubrica{Limite di risoluzione}
Distanza minima al di sotto della quale non siamo più in grado di vedere due punti distinti. Definisce l'ambito d'impiego dei tre strumenti:
\begin{itemize}
  \item \textbf{occhio umano}: strutture dal millimetro/centimetro in su;
  \item \textbf{microscopio ottico}: scala del micrometro;
  \item \textbf{microscopio elettronico}: scala del nanometro.
\end{itemize}

Dimensioni di riferimento:
\begin{table}[htbp]
  \centering \small \renewcommand{\arraystretch}{1.3}
  \begin{tabularx}{\textwidth}{|X|l|l|}
    \hline
    \textbf{Struttura} & \textbf{Dimensione} & \textbf{Strumento} \\
    \hline
    Uovo di gallina (tuorlo) & 3 cm & occhio nudo \\
    \hline
    Uova di rana o pesce & 1 mm & occhio nudo \\
    \hline
    Ovulo umano & 100\,$\mu$m & microscopio ottico \\
    \hline
    Tipica cellula vegetale & 10--100\,$\mu$m & microscopio ottico \\
    \hline
    Tipica cellula animale & 5--30\,$\mu$m & microscopio ottico \\
    \hline
    Cloroplasto & 2--10\,$\mu$m & microscopio ottico \\
    \hline
    Mitocondrio & 1--5\,$\mu$m & microscopio ottico \\
    \hline
    \textit{Escherichia coli} & 1\,$\mu$m & microscopio ottico \\
    \hline
    Grandi virus (HIV, influenza) & 100 nm & microscopio elettronico \\
    \hline
    Ribosoma & 25 nm & microscopio elettronico \\
    \hline
    Membrana cellulare (spessore) & 7--10 nm & microscopio elettronico \\
    \hline
    DNA a doppia elica (diametro) & 2 nm & microscopio elettronico \\
    \hline
    Atomo d'idrogeno & 0,1 nm & microscopio elettronico \\
    \hline
  \end{tabularx}
\end{table}

\rubrica{Microscopio ottico}
Componenti (dal basso): sorgente luminosa $\rightarrow$ condensatore $\rightarrow$ campione $\rightarrow$ lente dell'obiettivo $\rightarrow$ lente dell'oculare. Permette di vedere cellule colorate o vive, ma a risoluzione relativamente bassa.

Tecniche di microscopia ottica:
\begin{itemize}
  \item \textbf{campo luminoso}: la luce passa direttamente attraverso il campione; molte strutture hanno contrasto troppo basso $\rightarrow$ si usa un colorante, che però di norma fissa e uccide le cellule;
  \item \textbf{campo scuro}: il campione è illuminato con una inclinazione specifica e solo la luce da esso diffusa raggiunge la lente $\rightarrow$ immagine luminosa su sfondo nero;
  \item \textbf{contrasto di fase}: le differenze di rifrazione dovute alla densità del campione sono rese come differenze di contrasto $\rightarrow$ evidenzia strutture altrimenti invisibili; le cellule vive possono essere fotografate o filmate;
  \item \textbf{Nomarski} (contrasto a interferenza differenziale): lenti speciali amplificano le differenze di densità $\rightarrow$ aspetto tridimensionale;
  \item \textbf{fluorescenza}: strutture o molecole sono marcate con coloranti fluorescenti che emettono luce sotto illuminazione ultravioletta $\rightarrow$ localizzazione delle molecole marcate;
  \item \textbf{confocale a scansione laser}: un laser scandisce il campione marcato e il computer focalizza la luce su un singolo piano $\rightarrow$ immagine tridimensionale più nitida delle altre tecniche ottiche.
\end{itemize}

\rubrica{Microscopio elettronico}
\begin{itemize}
  \item \textbf{a trasmissione (MET)}: il fascio elettronico attraversa il campione $\rightarrow$ immagine ad altissima risoluzione, fortemente ingrandita (solo una parte di una sezione sottile);
  \item \textbf{a scansione (MES)}: sfrutta gli elettroni secondari $\rightarrow$ chiara visione delle caratteristiche della superficie della cellula.
\end{itemize}
A parità di ingrandimento (4500 volte), il microscopio elettronico mostra molti più dettagli del microscopio ottico (miofibrille, mitocondri, membrane delle singole strutture): il microscopio ottico non può fornire informazioni migliori.

\section{La materia vivente}

% ---------------------------------------------------------------------------
\subsection{Proprietà e componenti}

\rubrica{Proprietà della materia vivente}
\begin{itemize}
  \item metabolismo;
  \item accrescimento;
  \item moltiplicazione;
  \item irritabilità.
\end{itemize}

\textbf{Cellula}: la più piccola parte di protoplasma capace di un'esistenza indipendente; è l'unità morfologica, funzionale e fisiologica di tutti gli organismi viventi.

\rubrica{Componenti}
Il 75--85\% del protoplasma è costituito da acqua.
\begin{itemize}
  \item \textbf{inorganici}:
    \begin{itemize}
      \item elementi primari (97--98\%): C, H, O, N, S, P;
      \item cationi: $K^+$, $Na^+$, $Ca^{2+}$, $Mg^{2+}$;
      \item anioni: $Cl^-$, $HCO_3^-$, $PO_4^{3-}$, $SO_4^{3-}$;
      \item oligoelementi: Mn, Fe, Co, Cu, Mo, I, Br, V, Al, Se;
    \end{itemize}
  \item \textbf{organici}: carboidrati, lipidi, proteine, acidi nucleici.
\end{itemize}

% ---------------------------------------------------------------------------
\subsection{Glicidi o carboidrati}

Costituiscono più della metà dell'apporto energetico totale; sono formati da carbonio, idrogeno e ossigeno. Si dividono in monosaccaridi, disaccaridi/oligosaccaridi e polisaccaridi.

\rubrica{Monosaccaridi}
\begin{itemize}
  \item \textbf{triosi}: gliceraldeide;
  \item \textbf{pentosi}: ribosio;
  \item \textbf{esosi}: glucosio.
\end{itemize}
Note:
\begin{itemize}
  \item apporto giornaliero di carboidrati: 250--800 g nella dieta dei paesi occidentali;
  \item il glucosio è il nutriente fondamentale;
  \item esosi e pentosi: ruolo fondamentale nei processi nutritivi cellulari;
  \item glucosio e fruttosio: presenti in forma libera in alcuni frutti e vegetali.
\end{itemize}

\rubrica{Disaccaridi e oligosaccaridi}
Associazione di 2 o più molecole di monosaccaridi:
\begin{itemize}
  \item \textbf{saccarosio}: glucosio $+$ fruttosio;
  \item \textbf{maltosio}: glucosio $+$ glucosio;
  \item \textbf{lattosio}: glucosio $+$ galattosio.
\end{itemize}

\rubrica{Polisaccaridi}
Concatenazione di numerose unità monosaccaridiche:
\begin{itemize}
  \item \textbf{omopolisaccaridi}: unità monosaccaridiche identiche;
  \item \textbf{eteropolisaccaridi}: unità monosaccaridiche diverse.
\end{itemize}

\textbf{Amido}: principale polisaccaride di deposito nei vegetali. \textbf{Glicogeno}: principale polisaccaride di deposito negli animali. Entrambi costituiti da amilosio e amilopectina:
\begin{itemize}
  \item \textit{amilosio}: lineare, molecole di glucosio unite da legame monoglicosidico $\alpha(1\rightarrow4)$;
  \item \textit{amilopectina}: ramificata, legami glicosidici $\alpha(1\rightarrow6)$ nei punti di ramificazione.
\end{itemize}

\textbf{Cellulosa}: polimero lineare formato da migliaia di monomeri di glucosio uniti da legame $\beta(1\rightarrow4)$.
\begin{itemize}
  \item le catene si dispongono parallelamente mediante legami idrogeno, formando fibrille;
  \item funzione strutturale nel regno vegetale: costituisce la parete primaria e secondaria delle cellule vegetali;
  \item fermentata dalla flora batterica intestinale.
\end{itemize}

% ---------------------------------------------------------------------------
\subsection{Lipidi}

Due grandi classi:
\begin{itemize}
  \item \textbf{triacilgliceroli} (trigliceridi) --- 98\%: funzione di riserva energetica; lipidi di deposito, negli adipociti del tessuto adiposo; grassi animali;
  \item \textbf{fosfolipidi, glicolipidi, colesterolo, vitamine liposolubili} --- 2\%: funzione strutturale (membrane cellulari); dal colesterolo derivano la sintesi dei sali biliari e della vitamina D.
\end{itemize}

\rubrica{Acidi grassi}
Catene lineari con numero \textbf{pari} di atomi di carbonio; presentano un gruppo carbossilico (--COOH).

\textit{Saturi}:
\begin{itemize}
  \item catena corta $\rightarrow$ solubili in acqua e volatili: formico ($HCOOH$), acetico ($CH_3COOH$), propionico ($CH_3CH_2COOH$), butirrico ($CH_3CH_2CH_2COOH$);
  \item catena media: da 6 a 12 C (caprilico, laurico);
  \item catena lunga $\rightarrow$ insolubili in acqua: da 14 a 18 C (palmitico, stearico);
  \item catena molto lunga: da 20 a 24 C (arachidico).
\end{itemize}

\textit{Insaturi}: presenza di uno o più doppi legami; la solubilità in acqua aumenta con l'aumentare dei doppi legami:
\begin{itemize}
  \item monoinsaturo: acido oleico;
  \item polinsaturo: acido linoleico.
\end{itemize}

\rubrica{Lipidi semplici}
Derivano dall'esterificazione degli acidi grassi con alcoli:
\[
\text{glicerolo} + \text{acido grasso} \rightarrow \text{monogliceride} + \text{acqua}.
\]
\textbf{Sintesi dei trigliceridi}: glicerolo $+$ 3 acidi grassi $\rightarrow$ trigliceride $+$ 3 $H_2O$.

\textbf{Triacilgliceroli}: presenti nel tessuto adiposo; funzioni: riserva energetica, isolamento termico, sostegno di organi.

\rubrica{Lipidi complessi}
\textbf{Fosfolipidi}: molecole anfipatiche:
\begin{itemize}
  \item testa \textbf{polare}: gruppo fosfato (con glicerolo);
  \item code \textbf{apolari}: catene di acidi grassi.
\end{itemize}
Esempi: fosfatidilcolina (PC), fosfatidilinositolo (PI), fosfatidilserina (PS), fosfatidiletanolammina (PE), sfingomielina (SM).

\rubrica{Steroidi}
Lipidi policiclici la cui struttura di base è data da quattro anelli fusi tra loro. Esempi: colesterolo, vitamina D.

% ---------------------------------------------------------------------------
\subsection{Amminoacidi}

Struttura generale: un carbonio centrale ($\alpha$) legato a quattro gruppi:
\begin{itemize}
  \item gruppo carbossilico (--COOH);
  \item gruppo amminico (--$NH_2$);
  \item gruppo laterale R, proprio di ciascun amminoacido;
  \item un atomo di idrogeno.
\end{itemize}
La natura delle catene laterali conferisce proprietà diverse alle proteine in cui sono presenti.

\rubrica{Classificazione in base alla catena laterale}
\begin{itemize}
  \item \textbf{non polari}: alanina, valina, leucina, metionina, isoleucina, prolina, fenilalanina, triptofano;
  \item \textbf{polari ma non carichi}: cisteina, glicina, serina, treonina, asparagina, glutammina, tirosina;
  \item \textbf{carichi negativamente}: glutammato, aspartato;
  \item \textbf{carichi positivamente}: arginina, lisina, istidina.
\end{itemize}

\rubrica{Classificazione nutrizionale}
\begin{itemize}
  \item \textbf{essenziali} (l'organismo umano non è in grado di sintetizzarli, essenziali in ogni circostanza): fenilalanina, isoleucina, leucina, lisina, metionina, triptofano, treonina, valina, istidina;
  \item \textbf{condizionatamente essenziali} (richiedono un apporto esogeno, la sintesi endogena non è sufficiente): arginina, cisteina, glutammina, glicina, prolina, tirosina;
  \item \textbf{non essenziali} (sintetizzati dall'organismo in quantità sufficienti): alanina, aspartato, asparagina, glutammato, serina.
\end{itemize}

% ---------------------------------------------------------------------------
\subsection{Peptidi e legame peptidico}

\textbf{Legame peptidico}: si forma dalla condensazione del gruppo carbossilico (--COOH) del primo amminoacido con il gruppo amminico (--$NH_2$) del secondo, con formazione del legame --CO--NH-- e rilascio di una molecola d'acqua (reazione di sintesi disidratativa).
\begin{itemize}
  \item estremità \textbf{N-terminale}: gruppo amminico libero;
  \item estremità \textbf{C-terminale}: gruppo carbossilico libero.
\end{itemize}

% ---------------------------------------------------------------------------
\subsection{Proteine}

Quattro livelli di organizzazione strutturale.

\rubrica{Struttura primaria}
Data dalla sequenza amminoacidica; specifica di ogni proteina.

\rubrica{Struttura secondaria}
Ripiegamenti regolari della catena, mantenuti da legami idrogeno:
\begin{itemize}
  \item \textbf{$\alpha$-elica}: avvolgimento a spirale; i legami idrogeno mantengono l'avvolgimento; i gruppi laterali si proiettano all'esterno dai lati;
  \item \textbf{foglietto $\beta$}: la catena polipeptidica si ripiega su se stessa; metà dei gruppi laterali proiettata al di sopra del foglietto, l'altra metà al di sotto; i filamenti possono disporsi in direzioni opposte (antiparallelo) o nella stessa direzione (parallelo).
\end{itemize}

\rubrica{Struttura terziaria}
Organizzazione tridimensionale complessiva; dipende dalle interazioni tra i gruppi laterali:
\begin{itemize}
  \item legame a idrogeno;
  \item legame ionico;
  \item ponte disolfuro (S--S);
  \item interazione idrofobica.
\end{itemize}
Con la \textbf{denaturazione} la proteina perde la propria organizzazione ($\alpha$-eliche e filamenti $\beta$); durante la \textbf{rinaturazione} i ponti disolfuro aiutano a recuperare la forma nativa.

\rubrica{Struttura quaternaria}
Formata dall'unione di più catene polipeptidiche. Esempi: emoglobina, collagene.

\rubrica{Funzioni delle proteine}
\begin{itemize}
  \item \textbf{ormoni}: veicolano segnali regolatori tra le cellule;
  \item \textbf{anticorpi}: difesa;
  \item \textbf{accumulo}: riserva di molecole;
  \item \textbf{strutturali}: sostegno;
  \item \textbf{enzimatiche}: catalizzatori;
  \item \textbf{trasporto}: proteine canale;
  \item \textbf{motrici}: movimento cellulare;
  \item \textbf{regolatori}: regolano l'attività di altre molecole;
  \item \textbf{recettori}: ricevono segnali dall'ambiente extracellulare.
\end{itemize}

% ---------------------------------------------------------------------------
\subsection{Acidi nucleici}

Due tipi:
\begin{itemize}
  \item \textbf{DNA}: acido desossiribonucleico;
  \item \textbf{RNA}: acido ribonucleico.
\end{itemize}

\rubrica{Nucleosidi e nucleotidi}
\begin{itemize}
  \item \textbf{nucleoside}: zucchero (pentoso) $+$ base azotata, uniti da legame glicosidico;
  \item il pentoso è \textit{ribosio} (--OH in $2'$) o \textit{desossiribosio} (--H in $2'$);
  \item \textbf{nucleotide}: nucleoside $+$ gruppo fosfato $\rightarrow$ nucleoside monofosfato (NMF), difosfato (NDF), trifosfato (NTF).
\end{itemize}

\rubrica{Basi azotate}
\begin{itemize}
  \item \textbf{puriniche}: adenina, guanina;
  \item \textbf{pirimidiniche}: citosina, uracile, timina.
\end{itemize}
Appaiamento tra basi complementari mediante legami idrogeno: A--T e G--C nel DNA; A--U nell'RNA.

% ---------------------------------------------------------------------------
\subsection{Il DNA}

\rubrica{Struttura}
\begin{itemize}
  \item doppia elica; scheletro fosfato--desossiribosio, con legame fosfodiesterico;
  \item i due filamenti sono antiparalleli ($5'\rightarrow3'$);
  \item larghezza 2,2 nm; passo dell'elica $=$ 10 paia di basi $=$ 3,4 nm (20\,\AA{} di larghezza, 34\,\AA{} di passo, 3,4\,\AA{} per base).
\end{itemize}
Struttura proposta da Watson e Crick (\textit{Nature} 171, 737--738, 25 aprile 1953); l'analisi del DNA mediante diffrazione dei raggi X fu ottenuta da Rosalind Franklin.

\rubrica{Tre funzioni del materiale genetico}
\begin{enumerate}
  \item il DNA è depositario dell'informazione che codifica i caratteri ereditari;
  \item il DNA contiene l'informazione per dirigere la propria duplicazione;
  \item il DNA contiene l'informazione per dirigere la costruzione di proteine specifiche.
\end{enumerate}

\rubrica{Organizzazione della cromatina}
$DNA \rightarrow$ nucleosoma $\rightarrow$ fibra da 10 nm $\rightarrow$ fibra da 30 nm (solenoide) $\rightarrow$ domini ad ansa $\rightarrow$ impalcatura proteica $\rightarrow$ cromosoma metafasico.
\begin{itemize}
  \item \textbf{nucleosoma}: DNA avvolto attorno a un core di 8 istoni (2 molecole ciascuno di H2A, H2B, H3, H4);
  \item l'istone \textbf{H1} è associato al DNA linker tra i nucleosomi.
\end{itemize}

\rubrica{Replicazione}
\begin{enumerate}
  \item i due filamenti si svolgono e si separano;
  \item ciascun filamento funge da stampo per l'aggiunta di basi complementari (A--T, G--C);
  \item si formano 2 nuove doppie eliche identiche a quella parentale.
\end{enumerate}
Ogni molecola risultante è metà ``vecchia'' e metà ``nuova'' $\rightarrow$ replicazione \textbf{semiconservativa}.

\rubrica{DNA superavvolto}
Il DNA può presentarsi in forma rilassata o superavvolta (fortemente condensata). Esempio di compattazione: il DNA rilasciato dalla testa del batteriofago T2 è lungo 68\,$\mu$m, circa 60 volte la testa del fago nella quale è contenuto.

% ---------------------------------------------------------------------------
\subsection{L'RNA}

Caratteristiche:
\begin{itemize}
  \item singola catena nucleotidica;
  \item zucchero: ribosio;
  \item base: uracile (al posto della timina).
\end{itemize}

Diversi tipi:
\begin{itemize}
  \item \textbf{mRNA} (messaggero);
  \item \textbf{tRNA} (transfer): struttura a trifoglio; presenta l'\textit{anticodone} e, all'estremità $3'$, il sito di attacco dell'amminoacido;
  \item \textbf{rRNA} (ribosomiale): componente delle subunità ribosomiali.
\end{itemize}

\section{Gli organismi viventi}

Gli organismi viventi sono in grado di:
\begin{itemize}
  \item riprodursi e trasmettere i propri caratteri (eredità);
  \item compiere un ciclo vitale;
  \item rinnovare continuamente la propria struttura;
  \item reagire agli stimoli;
  \item muoversi;
  \item evolversi.
\end{itemize}

% ---------------------------------------------------------------------------
\subsection{Riproduzione}

Insieme dei meccanismi mediante i quali gli esseri viventi provvedono alla conservazione della propria specie, generando nuovi individui simili a sé che subentreranno al genitore, o ai genitori, nella popolazione.

\rubrica{Riproduzione asessuata}
Genera individui che mantengono invariato il patrimonio genetico del genitore.
\begin{itemize}
  \item \textbf{moltiplicazione cellulare} (divisione mitotica): una cellula madre si divide in due cellule figlie identiche; negli organismi unicellulari è l'unica forma di riproduzione e incrementa una popolazione geneticamente uguale, nei pluricellulari è il meccanismo dell'accrescimento del singolo individuo;
  \item \textbf{scissione}: un individuo generante, dividendo il proprio corpo, genera individui figli identici tra loro e al generante;
  \item \textbf{frammentazione}: una parte dell'organismo si stacca e forma un nuovo individuo completo e identico al generante (per rigenerazione);
  \item \textbf{gemmazione}: formazione di gemme laterali che si separano definitivamente mediante una strozzatura, dando origine a nuovi individui uguali al generante;
  \item \textbf{poliembrionia}: divisione in due o più parti dello zigote o dell'embrione nei primissimi stadi di sviluppo $\rightarrow$ gemelli monozigoti;
  \item \textbf{sporulazione}: cellule specializzate (sporocisti), in condizioni ambientali sfavorevoli, producono per divisione mitotica speciali cellule riproduttive (\textit{mitospore}) capaci di generare un nuovo individuo quando il contesto ridiventa favorevole; le mitospore hanno una spessa parete protettiva che permette di resistere alle condizioni avverse.
\end{itemize}

\rubrica{Riproduzione sessuata}
Il nuovo individuo si origina dalla fusione dei gameti; la divisione che produce i gameti è la \textbf{meiosi}. Aumenta la variabilità genetica all'interno di una specie.
\begin{itemize}
  \item \textbf{anfigonia}: fecondazione di due gameti per formare una nuova cellula, lo zigote; la fusione di due gameti aploidi genera una cellula diploide (lo zigote), con metà del materiale genetico dal gamete maschile e metà da quello femminile;
  \item \textbf{partenogenesi}: riproduzione tramite sviluppo di gameti femminili in assenza di fecondazione; avviene generalmente per auto-attivazione dell'uovo non fecondato che, per restituzione anafasica, ricostituisce un genoma diploide omozigote in tutti i loci (due assetti aploidi identici). Generalmente è un evento accidentale, ma in alcune specie è una forma alternativa all'anfigonia e genera prole: es.\ \textit{Bacillus rossius} (insetto stecco) e alcuni imenotteri sociali (api, vespe).
\end{itemize}

% ---------------------------------------------------------------------------
\subsection{Ciclo vitale}

$\text{nascita} \rightarrow \text{crescita} \rightarrow \text{sviluppo} \rightarrow \text{morte}$; durata variabile a seconda della specie.
\begin{itemize}
  \item \textbf{crescita}: aumento della dimensione cellulare e/o del numero di cellule;
  \item \textbf{sviluppo}: modificazioni strutturali e funzionali, fino a un progressivo rallentamento e deterioramento delle funzionalità che porta alla morte.
\end{itemize}

% ---------------------------------------------------------------------------
\subsection{Metabolismo}

Il complesso delle trasformazioni di natura chimica che avvengono negli organismi viventi.
\begin{itemize}
  \item \textbf{capacità di nutrirsi}: organismi autotrofi, organismi eterotrofi;
  \item \textbf{capacità di trasformare materia ed energia} (insieme delle reazioni chimiche e fisiche che avvengono in un organismo o in una sua parte):
    \begin{itemize}
      \item \textit{anabolismo}: produce molecole complesse a partire da molecole più semplici, con consumo di energia (reazioni endoergoniche); es.\ biosintesi di una proteina;
      \item \textit{catabolismo}: degrada molecole complesse in molecole più semplici (reazioni esoergoniche), liberando energia; es.\ respirazione cellulare o ciclo di Krebs.
    \end{itemize}
\end{itemize}

% ---------------------------------------------------------------------------
\subsection{Reazione agli stimoli}

Capacità di percepire e reagire a stimoli esterni: cambiamenti di temperatura e di pressione, presenza/assenza di luce, cambiamenti chimici. Esempi: formazione di spore, letargo, mimetismo, reazione ai rumori o a stimoli olfattivi, fototropismo.

% ---------------------------------------------------------------------------
\subsection{Movimento}

\begin{itemize}
  \item \textbf{organismi unicellulari}: movimento ameboide (caratteristico delle \textit{Amoebe}, che si muovono emettendo prolungamenti detti pseudopodi); oppure mediante ciglia o flagelli;
  \item \textbf{organismi superiori}:
    \begin{itemize}
      \item \textit{piante}: allungamenti o rotazioni; movimenti lenti e limitati, in funzione della presenza di acqua (radici) o del sole (fronde, fiori);
      \item \textit{animali}: organi preposti al movimento, che può essere rapidissimo e coprire notevoli distanze (pinne, ali, zampe, arti).
    \end{itemize}
\end{itemize}

% ---------------------------------------------------------------------------
\subsection{Ambienti biologici}

\begin{itemize}
  \item \textbf{ambiente acqueo}: marino, d'acqua dolce;
  \item \textbf{ambiente terrestre}: epigeo, ipogeo.
\end{itemize}

\rubrica{Ambiente marino}
\begin{itemize}
  \item \textbf{dominio bentonico} (o di fondo): ambiente dove vivono gli organismi legati più o meno direttamente ai fondali:
    \begin{itemize}
      \item \textit{ambiente litoraneo}: acque e fondali che circondano le coste (la cosiddetta \textit{platea continentale}), fino a una profondità massima di 200 m; molto eterogeneo (in superficie risente dei moti ondosi ed è ben illuminato, in profondità è quasi buio e il moto ondoso quasi assente); i livelli meno profondi sono i più ricchi di vita animale e vegetale;
      \item \textit{ambiente profondo}: oltre la platea continentale il mare sprofonda più rapidamente (a gradino);
    \end{itemize}
  \item \textbf{dominio pelagico}: acque libere, distanti dalle coste e dal fondo, che si estendono dalla superficie fino agli abissi delle fosse oceaniche; vi vivono gli organismi che conducono una vita non vincolata in maniera esclusiva al fondale.
\end{itemize}
Nell'ambiente marino la variabilità maggiore è data da salinità e ossigenazione.

\rubrica{Organismi dell'ambiente marino}
\begin{itemize}
  \item \textbf{plancton}: organismi animali e vegetali che galleggiano, trasportati quasi tutti passivamente da moti ondosi e correnti (alghe unicellulari, protozoi, larve, piccoli crostacei, meduse, alghe pluricellulari);
  \item \textbf{necton}: organismi che nuotano attivamente, provvisti di mezzi per nuotare, emergere e immergersi (pesci, cefalopodi, tartarughe, cetacei);
  \item \textbf{benthos}: organismi che vivono in stretto contatto con il fondo o fissati a un substrato solido (alghe pluricellulari, animali che camminano o strisciano come i crostacei, animali sessili ancorati a un substrato e incapaci di muoversi come i molluschi bivalvi).
\end{itemize}

\rubrica{Ambienti d'acqua dolce}
\begin{itemize}
  \item \textbf{laghi}: ambiente in parte sovrapponibile a quello marino; i laghi di montagna, più freddi, hanno flora e fauna limitata rispetto a quelli di pianura;
  \item \textbf{fiumi, torrenti e sorgenti}: flora e fauna diverse a seconda del corso (lento o impetuoso) e delle acque (calde o fredde).
\end{itemize}

\rubrica{Ambiente terrestre}
Varia in base a latitudine, longitudine, sottosuolo, vegetazione, clima e stagioni; variabilità maggiore data da temperatura e umidità.
\begin{itemize}
  \item \textbf{epigeo}: organismi che vivono sulla superficie della terra;
  \item \textbf{ipogeo}: organismi che vivono dentro la terra (grotte, anfratti, gallerie scavate nel sottosuolo);
  \item \textbf{entozoico}: endoparassiti che vivono all'interno di un altro animale ospite.
\end{itemize}

% ---------------------------------------------------------------------------
\subsection{Rapporti intraspecifici}

Rapporti che intercorrono tra individui della stessa specie.
\begin{itemize}
  \item \textbf{colonie}: individui della stessa specie che, dopo essersi riprodotti per via asessuata (gemmazione), non si separano e rimangono fisicamente uniti:
    \begin{itemize}
      \item \textit{colonie omeomorfe}: tutti gli individui sono uguali (spugne, madrepore, coralli);
      \item \textit{colonie eteromorfe}: gruppi di individui si specializzano in una particolare funzione fino a una spiccata differenziazione morfo-funzionale (``superorganismi'');
    \end{itemize}
  \item \textbf{società animali}: rapporti tra animali della stessa specie che, pur fisicamente separati, conducono vita comune (temporanee, durature, individualiste, collettiviste).
\end{itemize}

% ---------------------------------------------------------------------------
\subsection{Rapporti interspecifici}

Rapporti tra individui di specie diversa che si instaurano all'interno di uno stesso ambiente; riguardano la nutrizione e l'occupazione degli spazi.

\rubrica{Simbiosi}
Stretta relazione fra individui di specie diverse che convivono per trarre un beneficio reciproco.
\begin{itemize}
  \item \textbf{mutualismo}: nessuno degli individui è in grado di vivere isolato (licheni; paguro e attinia; \textit{Rhizobium});
  \item \textbf{commensalismo}: solo una delle due specie trae vantaggio dall'associazione, senza però danneggiare l'altra (pesce pagliaccio; acari foretici; storno);
  \item \textbf{inquilinismo}: i rapporti fra le due specie si riducono all'occupazione di uno spazio comune.
\end{itemize}

\rubrica{Amensalismo}
Una specie impedisce o diminuisce il successo di un'altra, senza però trarne vantaggio. Si ha quando un organismo produce una sostanza chimica, come parte del suo normale metabolismo, con effetto negativo su altri organismi (es.\ \textit{Penicillium}, che secerne penicillina, un potente battericida).

\rubrica{Parassitismo}
Una delle due specie conviventi, il parassita, trae vantaggio dall'associazione a scapito dell'altra, l'ospite, creandole un danno biologico.
\begin{itemize}
  \item il parassita dipende dall'ospite, a cui è più o meno intimamente legato da una relazione anatomica e fisiologica;
  \item il parassita ha una struttura anatomica più semplice di quella dell'ospite;
  \item l'ospite è più grande del parassita;
  \item il ciclo vitale del parassita è più breve di quello dell'ospite e si conclude prima della sua morte;
  \item il parassita ha rapporti con un solo ospite, ma non viceversa.
\end{itemize}
Tipi:
\begin{itemize}
  \item \textbf{facoltativo}: non si hanno modificazioni morfo-funzionali;
  \item \textbf{obbligato}: la vita del parassita è subordinata a quella dell'ospite;
  \item \textbf{ectoparassiti}: adottano elaborati meccanismi e strategie per attaccare un ospite (es.\ le sanguisughe individuano l'ospite tramite sensori di movimento e ne accertano l'identità attraverso la temperatura della pelle e indicazioni chimiche; hanno un adattamento morfofunzionale per aderire all'ospite e succhiarne il nutrimento);
  \item \textbf{endoparassiti}: attaccano l'ospite in modo passivo (es.\ gli endoparassiti dell'intestino umano infestano l'uomo tramite l'ingestione di acqua sporca o cibi infestati; hanno subìto grosse modificazioni morfofunzionali).
\end{itemize}
Esempi: \textit{Taenia solium}, \textit{Plasmodium malariae}.

\section{Citologia}

Schema generale dei tipi di cellula: cellula batterica, cellula vegetale, cellula animale. Oltre alle cellule esistono anche batteriofagi, virus, viroidi e prioni.

% ---------------------------------------------------------------------------
\subsection{Virus e batteriofagi}

\begin{itemize}
  \item dimensioni: 10--300 nm;
  \item \textbf{virus}: si riproducono all'interno di cellule eucariotiche;
  \item \textbf{batteriofagi} (o fagi): si riproducono all'interno dei batteri;
  \item utilizzano l'apparato sintetico della cellula ospite per sintetizzare e assemblare le molecole che li compongono.
\end{itemize}

\rubrica{Struttura del virus}
\begin{itemize}
  \item \textbf{genoma virale}: DNA o RNA;
  \item \textbf{capside}: involucro proteico costituito da subunità (capsomeri) disposte in modo regolare; protegge il genoma virale in ambiente extracellulare e consente la penetrazione del virus nella cellula ospite;
  \item \textbf{pericapside} (o peplos): ulteriore involucro lipoproteico, talvolta presente all'esterno del capside.
\end{itemize}

\rubrica{Classificazione dei virus}
\begin{itemize}
  \item in base al tipo di cellula infettata: batteriofagi/fagi, virus vegetali, virus animali;
  \item in base all'acido nucleico: desossiribovirus (virus a DNA), ribovirus o retrovirus (virus a RNA);
  \item in base alla forma: allungati, sferici, poliedrici;
  \item in base al tipo di simmetria: cubica, elicoidale, complessa.
\end{itemize}

% ---------------------------------------------------------------------------
\subsection{Replicazione di virus e batteriofagi}

\rubrica{Ciclo litico dei batteriofagi}
\begin{enumerate}
  \item aggancio: il fago aderisce alla superficie cellulare del batterio;
  \item penetrazione: il DNA del fago entra nella cellula batterica;
  \item replicazione e sintesi: il DNA del fago viene replicato e le sue proteine sintetizzate;
  \item assemblaggio: le componenti del fago vengono assemblate in nuovi virus;
  \item rilascio: la cellula batterica si lisa liberando molti fagi, che possono infettare altre cellule.
\end{enumerate}

\rubrica{Ciclo lisogenico dei batteriofagi}
\begin{enumerate}
  \item aggancio;
  \item penetrazione;
  \item integrazione: il DNA fagico si integra nel DNA batterico (\textit{profago});
  \item replicazione: il profago integrato si replica quando viene replicato il DNA batterico; le cellule figlie possono esibire nuove proprietà.
\end{enumerate}
Il fago \textbf{lambda} è un esempio di fago temperato, che può effettuare il ciclo litico o quello lisogenico.

\rubrica{Replicazione a confronto}
Nel \textbf{virus}: adesione a parete o membrana cellulare $\rightarrow$ rilascio di enzimi che consentono l'adesione alla cellula ospite $\rightarrow$ dissoluzione del capside $\rightarrow$ ingresso del genoma virale $\rightarrow$ duplicazione del genoma $\rightarrow$ sintesi delle proteine del capside $\rightarrow$ assemblaggio delle nuove particelle virali.

Nel \textbf{batteriofago}: adesione alla parete batterica $\rightarrow$ iniezione dell'acido nucleico nel batterio $\rightarrow$ duplicazione del genoma $\rightarrow$ sintesi delle proteine del capside $\rightarrow$ assemblaggio dei nuovi fagi $\rightarrow$ lisi della cellula ospite $\rightarrow$ fuoriuscita dei fagi, che possono infettare altri batteri.

% ---------------------------------------------------------------------------
\subsection{Viroidi e prioni}

\rubrica{Viroidi}
\begin{itemize}
  \item piccole molecole di RNA chiuse ad anello, capaci di autoreplicazione;
  \item un solo acido nucleico;
  \item privi di capside;
  \item infettano preferenzialmente le cellule vegetali.
\end{itemize}

\rubrica{Prioni}
\begin{itemize}
  \item molecole proteiche normalmente presenti in soggetti sani;
  \item in particolari condizioni alterano la propria configurazione spaziale;
  \item trasmettono la propria configurazione alterata alle molecole sane.
\end{itemize}

% ---------------------------------------------------------------------------
\subsection{Cellula procariote}

Componenti: parete cellulare, membrana plasmatica, capsula, DNA (nel nucleoide/area nucleare), plasmidi, ribosomi, citoplasma, mesosoma, pili, flagello, granuli di riserva.

\rubrica{Classificazione dei batteri per forma}
\begin{itemize}
  \item \textbf{bacilli}: a bastoncino;
  \item \textbf{cocchi}: a sfera;
  \item \textbf{diplococchi}: due cocchi uniti;
  \item \textbf{streptococchi}: disposti a catena;
  \item \textbf{stafilococchi}: disposti a grappolo;
  \item \textbf{spirilli}: a spirale;
  \item \textbf{vibrioni}: a virgola;
  \item \textbf{spirochete}: con più curve.
\end{itemize}

\rubrica{Altre classificazioni dei batteri}
\begin{itemize}
  \item in base alla temperatura: \textit{criofili/psicrofili} (\textit{Pseudomonas}, \textit{Lactobacillus}), \textit{mesofili} (\textit{Salmonella}, stafilococco, streptococco), \textit{termofili} (\textit{Thermus aquaticus});
  \item in base alla presenza/assenza di ossigeno: aerobi, anaerobi, aerobi/anaerobi facoltativi;
  \item in base al grado di acidità: acidofili, alcalofili;
  \item in base alla concentrazione salina: alofili (elevate concentrazioni saline).
\end{itemize}

\rubrica{Parete cellulare batterica}
Costituita da uno strato di peptidoglicano.
\begin{itemize}
  \item \textbf{Gram-positivi}: parete spessa (spesso strato di peptidoglicano); la capsula può essere presente;
  \item \textbf{Gram-negativi}: parete sottile, con una membrana esterna oltre allo strato di peptidoglicano.
\end{itemize}

\rubrica{Duplicazione dei procarioti}
Avviene per scissione binaria: il cromosoma (DNA), collegato al mesosoma, si duplica a partire da un'origine di replicazione (replicazione bidirezionale); le due origini migrano verso i poli della cellula; la membrana si accresce verso l'interno e si forma una nuova parete $\rightarrow$ due cellule figlie che si separano.

% ---------------------------------------------------------------------------
\subsection{Cellula eucariote}

Organuli: nucleo (involucro nucleare, nucleoplasma, nucleolo, cromatina), membrana plasmatica, reticolo endoplasmatico (rugoso e liscio), ribosomi, complesso di Golgi, lisosomi, perossisomi, mitocondri, citoscheletro (microfilamenti, microtubuli, filamenti intermedi), centrioli, vescicole, citosol. La cellula \textbf{vegetale} presenta in più parete cellulare, cloroplasti, vacuolo centrale e plasmodesmi.

\rubrica{Membrana plasmatica}
Doppio strato fosfolipidico (fosfolipidi con testa idrofilica e coda idrofobica); costituisce il limite esterno del citoplasma; presenta canali proteici di trasporto per le sostanze idrosolubili, che non possono attraversare la regione fosfolipidica.

\rubrica{Involucro nucleare}
Sistema di due membrane concentriche attraversate da pori nucleari, che facilitano il trasporto di molecole tra nucleo e citoplasma; sulla superficie dell'involucro (membrana esterna, affacciata al citoplasma) sono presenti ribosomi.

\rubrica{Reticolo endoplasmatico}
\begin{itemize}
  \item \textbf{RE rugoso}: con ribosomi sulla superficie (sede della sintesi proteica);
  \item \textbf{RE liscio}: privo di ribosomi.
\end{itemize}

\rubrica{Complesso di Golgi}
Insieme di cisterne; riceve vescicole originate dal RE (modificazione e smistamento delle proteine) ed emette vescicole in fase di gemmazione.

\rubrica{Mitocondri}
Membrana esterna e membrana interna (che si ripiega formando le creste); matrice interna e compartimento intermembrana; sede del metabolismo energetico.

\rubrica{Cloroplasti}
Membrana esterna e membrana interna; tilacoidi impilati in \textit{grana} (granum); stroma (fluido interno).

\rubrica{Componenti cellulari a confronto}
Presenza (X) dei principali componenti nei procarioti (Eubatteri, Archea) e negli eucarioti (Protisti, Funghi, Piante, Animali).
\begin{table}[htbp]
  \centering \footnotesize \renewcommand{\arraystretch}{1.15}
  \begin{tabular}{|l|c|c|c|c|c|c|}
    \hline
    \textbf{Componente} & \textbf{Eubatt.} & \textbf{Archea} & \textbf{Protisti} & \textbf{Funghi} & \textbf{Piante} & \textbf{Animali} \\
    \hline
    Nucleoide            & X & X &        &        &   &   \\ \hline
    Nucleo               &   &   & X      & X      & X & X \\ \hline
    Involucro nucleare   &   &   & X      & X      & X & X \\ \hline
    Nucleolo             &   &   & X      & X      & X & X \\ \hline
    Membrana plasmatica  & X & X & X      & X      & X & X \\ \hline
    Parete cellulare     & X & X & alcuni & X      & X &   \\ \hline
    Ribosomi             & X & X & X      & X      & X & X \\ \hline
    Reticolo endoplasm.  &   &   & X      & X      & X & X \\ \hline
    Complesso di Golgi   &   &   & X      & X      & X & X \\ \hline
    Lisosoma             &   &   & X      & alcuni & X & X \\ \hline
    Mitocondrio          &   &   & X      & X      & X & X \\ \hline
    Cloroplasto          &   &   & alcuni &        & X &   \\ \hline
    Vacuolo centrale     &   &   &        &        & X &   \\ \hline
    Microfilamenti       &   &   & X      & X      &   & X \\ \hline
    Microtubuli          &   &   & X      & X      & X & X \\ \hline
    Filamenti intermedi  &   &   & X      &        &   & X \\ \hline
  \end{tabular}
\end{table}

\section{Biotecnologie}

% ---------------------------------------------------------------------------
\subsection{Microscopia e preparazione dei campioni}

Strumenti di indagine: microscopio ottico e microscopio elettronico (a trasmissione, a scansione).

\rubrica{Preparazione dei campioni per il microscopio elettronico}
$\text{fissazione} \rightarrow \text{disidratazione (etanolo 70--100\%)} \rightarrow \text{inclusione in resina} \rightarrow \text{sezionamento} \rightarrow \text{colorazione}$.
\begin{itemize}
  \item fissazione (es.\ glutaraldeide, $OsO_4$) e disidratazione in etanolo a concentrazione crescente;
  \item infiltrazione e inclusione in resina, polimerizzata in un blocco solido;
  \item sezionamento con \textbf{ultramicrotomo} in sezioni di $\sim$100 nm;
  \item colorazione delle sezioni con metalli pesanti.
\end{itemize}

% ---------------------------------------------------------------------------
\subsection{Centrifugazione in gradiente di densità}

Purificazione di frazioni subcellulari: in un gradiente di densità (saccarosio) i vari organelli sedimentano finché raggiungono il livello a densità uguale alla loro, dove formano bande.
\begin{itemize}
  \item lisosomi: 1,12 g/ml;
  \item mitocondri: 1,18 g/ml;
  \item perossisomi: 1,23 g/ml.
\end{itemize}

\rubrica{Sedimentazione degli acidi nucleici}
\begin{itemize}
  \item separazione del DNA in base alle \textbf{dimensioni} (velocità di sedimentazione);
  \item separazione in base alla \textbf{composizione in basi} (centrifugazione all'equilibrio).
\end{itemize}

\rubrica{Gradiente di cloruro di cesio (CsCl)}
Centrifugazione all'equilibrio: separa il DNA in base alla densità.
\begin{itemize}
  \item DNA ricco in AT: $\sim$1,65 g/ml; DNA ricco in GC: $\sim$1,75 g/ml;
  \item separa il DNA plasmidico dal DNA cromosomico batterico (due bande);
  \item il DNA è reso visibile con \textbf{bromuro di etidio}, fluorescente ai raggi ultravioletti.
\end{itemize}

% ---------------------------------------------------------------------------
\subsection{Enzimi di restrizione}

\textbf{Endonucleasi di restrizione}: tagliano il DNA a livello di sequenze palindromiche di riconoscimento, producendo estremità complementari dette \textit{estremità coesive} (``sticky ends'').
\begin{itemize}
  \item esempi: EcoRI (da \textit{E.\ coli}), HindII e HindIII (da \textit{Haemophilus influenzae}), HpaII, AluI;
  \item la \textbf{metilazione} della sequenza di riconoscimento protegge il DNA dal taglio.
\end{itemize}

% ---------------------------------------------------------------------------
\subsection{Clonaggio e DNA ricombinante}

\rubrica{Vettori plasmidici}
Plasmidi (es.\ pUC118, pUC119) che contengono:
\begin{itemize}
  \item una \textbf{regione di policlonaggio} con i siti di taglio per gli enzimi di restrizione;
  \item il gene \textit{lacZ};
  \item un'origine di replicazione;
  \item un gene di resistenza all'ampicillina.
\end{itemize}

\rubrica{Formazione di DNA ricombinante}
Lo stesso enzima di restrizione taglia sia il plasmide sia il DNA da inserire $\rightarrow$ i frammenti hanno estremità adesive complementari $\rightarrow$ il frammento estraneo si inserisce nel plasmide e le interruzioni vengono saldate dalla \textbf{DNA ligasi}.

% ---------------------------------------------------------------------------
\subsection{\texorpdfstring{PCR (\textit{Polymerase Chain Reaction})}{PCR (Polymerase Chain Reaction)}}

\textbf{Scopo}: amplificare in provetta una sequenza bersaglio, cioè produrne un grande numero di copie, senza clonaggio.

\rubrica{Componenti della miscela di reazione}
\begin{enumerate}
  \item il DNA contenente la sequenza bersaglio da amplificare;
  \item una coppia di primer, ciascuno complementare a una delle estremità della sequenza bersaglio;
  \item i quattro nucleotidi trifosfato (dATP, dCTP, dGTP, dTTP);
  \item DNA polimerasi \textbf{termostabile} (isolata da un microrganismo termofilo, perché la PCR usa alte temperature in alcuni passaggi).
\end{enumerate}
La reazione è ciclica: a ogni ciclo il numero di copie raddoppia ($2 \rightarrow 4 \rightarrow 8 \rightarrow \dots$).

% ---------------------------------------------------------------------------
\subsection{Elettroforesi su gel di agarosio}

Separa molecole di DNA, RNA o proteine in base alla dimensione, alla carica o ad altre proprietà, applicando un campo elettrico a un gel.
\begin{enumerate}
  \item si prepara un gel di agarosio con pozzetti e lo si immerge in un tampone tra due elettrodi;
  \item si caricano nei pozzetti i campioni (es.\ prodotti di PCR) e un \textit{marker} (frammenti di dimensione nota);
  \item applicando la corrente, i frammenti (carichi negativamente) migrano verso il polo positivo; i più corti migrano più velocemente;
  \item i frammenti della stessa lunghezza formano bande, visualizzate con un colorante che lega il DNA e diventa fluorescente ai raggi UV.
\end{enumerate}

% ---------------------------------------------------------------------------
\subsection{Sequenziamento del DNA}

Metodo con terminatori di catena: si preparano quattro reazioni di sintesi, ciascuna contenente DNA polimerasi, primer, i quattro dNTP e uno dei quattro \textbf{didesossiribonucleosidi trifosfato} (ddNTP), marcato.
\begin{itemize}
  \item l'incorporazione di un ddNTP interrompe la sintesi $\rightarrow$ si generano frammenti di lunghezza diversa;
  \item i prodotti sono separati per elettroforesi (gel di poliacrilammide o elettroforesi capillare) in base alla grandezza;
  \item la lettura fornisce la sequenza; nei sistemi automatici i ddNTP sono marcati con coloranti fluorescenti letti da un laser e un rilevatore collegato a un computer.
\end{itemize}

% ---------------------------------------------------------------------------
\subsection{Progetto Genoma Umano}

\begin{itemize}
  \item tempo previsto: 1990--2005; tempo effettivamente impiegato: 1990--2002;
  \item sequenziato tutto il genoma umano ($3 \cdot 10^9$ bp);
  \item solo il 3\% è costituito da DNA codificante;
  \item il restante 97\% è detto intergenico o non codificante.
\end{itemize}

% ---------------------------------------------------------------------------
\subsection{Organismi transgenici e clonazione}

\rubrica{Topi transgenici}
Scopo: produrre un animale transgenico che possa trasmettere il transgene alla progenie.
\begin{enumerate}
  \item si introduce il gene desiderato nelle cellule della linea germinale (prelevate da un embrione) tramite iniezione o elettroporazione;
  \item si clona la cellula che ha incorporato il transgene, ottenendo una coltura pura di cellule transgeniche;
  \item si iniettano le cellule transgeniche in embrioni precoci (blastocisti);
  \item si impiantano gli embrioni in madri adottive e si incrocia tra loro la progenie $\rightarrow$ animale ingegnerizzato geneticamente.
\end{enumerate}
Esempio: topo gigante ottenuto introducendo il gene di ratto per l'ormone della crescita.

\rubrica{Clonazione}
Febbraio 1997: clonazione della pecora \textbf{Dolly}, primo mammifero clonato a partire da una cellula di un adulto.

% ---------------------------------------------------------------------------
\subsection{Cellule staminali}

\begin{itemize}
  \item \textbf{cellule totipotenti} --- cellule staminali embrionali (CSE): in grado di produrre qualsiasi cellula dell'organismo;
  \item \textbf{cellule pluripotenti o multipotenti} --- cellule staminali dell'adulto (CSA): possono dare origine a un certo numero di tessuti (es.\ le staminali ematopoietiche del midollo osseo generano le cellule del sangue).
\end{itemize}


% =============================================================================
% CAPITOLO 2 — Biologia cellulare
% =============================================================================
\chapter{Biologia cellulare}
\section{La membrana plasmatica: struttura}

% ---------------------------------------------------------------------------
\subsection{Aspetto e composizione}

Al microscopio elettronico la membrana appare \textbf{trilaminare} (a tre strati), dopo colorazione con osmio: l'osmio si lega preferenzialmente ai gruppi delle teste polari del doppio strato lipidico.

Componenti:
\begin{itemize}
  \item \textbf{fosfolipidi} (testa idrofilica, coda idrofobica);
  \item \textbf{colesterolo};
  \item \textbf{glicolipidi};
  \item \textbf{proteine} associate alla membrana.
\end{itemize}

% ---------------------------------------------------------------------------
\subsection{Il doppio strato fosfolipidico}

Il \textbf{fosfolipide} è una molecola anfipatica, di forma pressoché cilindrica:
\begin{itemize}
  \item \textbf{testa idrofilica}: alcol polare (es.\ colina) $+$ gruppo fosfato $+$ glicerolo;
  \item \textbf{coda idrofobica}: catene di acidi grassi.
\end{itemize}
In acqua i fosfolipidi si associano spontaneamente in un \textbf{doppio strato}: le teste idrofiliche sono rivolte verso l'acqua, le code idrofobiche verso l'interno, non a contatto con l'acqua.
\begin{itemize}
  \item confronto: i \textit{detergenti}, anfipatici ma di forma conica, in acqua si associano formando strutture sferiche (micelle);
  \item il doppio strato può ripiegarsi e chiudersi a formare una \textbf{vescicola}.
\end{itemize}

% ---------------------------------------------------------------------------
\subsection{Modelli di membrana}

\begin{itemize}
  \item \textbf{modello di Davson-Danielli} (1954): entrambe le superfici del doppio strato lipidico sono rivestite da uno strato monomolecolare di proteine, che si estendono attraverso la membrana formando canali delimitati da proteine;
  \item \textbf{modello a mosaico fluido} (Singer e Nicolson, 1972): le proteine penetrano nel doppio strato lipidico.
\end{itemize}

% ---------------------------------------------------------------------------
\subsection{Le proteine di membrana}

Tre classi:
\begin{itemize}
  \item \textbf{integrali}: attraversano il doppio strato con una o più eliche transmembrana (\textit{single pass} = una sola elica; \textit{multiple pass} = più eliche); motivi strutturali: singola $\alpha$-elica, più $\alpha$-eliche, foglietto $\beta$ ripiegato a barilotto;
  \item \textbf{periferiche}: legate con legami non covalenti ai gruppi polari delle teste lipidiche e/o a una proteina integrale;
  \item \textbf{legate ai lipidi}: unite con legami covalenti a un gruppo lipidico inserito in uno dei due foglietti (es.\ ancoraggio GPI, acido grasso, gruppo prenile).
\end{itemize}

\rubrica{Struttura}
Le regioni ad $\alpha$-elica attraversano la membrana; la proteina presenta superfici \textbf{idrofobiche} (a contatto con le code lipidiche) e \textbf{idrofiliche} (verso l'acqua o verso un canale idrofilico che attraversa la membrana). Esempi: batteriorodopsina (assorbe l'energia luminosa negli archeobatteri fotosintetici); acquaporina (proteina formata da quattro subunità che circonda un canale acquoso, avvolta da uno strato di molecole lipidiche molto aderenti).

% ---------------------------------------------------------------------------
\subsection{Il colesterolo nella membrana}

Il gruppo --OH idrofilico si inserisce nella regione polare del doppio strato, mentre la struttura ciclica (idrofobica) si estende nell'interno apolare della membrana.

% ---------------------------------------------------------------------------
\subsection{Fluidità della membrana}

\rubrica{Esperimento di Frye-Edidin}
Fusione di una cellula umana (proteine di membrana marcate con colorante fluorescente rosso) e di una cellula di topo (marcate in verde) $\rightarrow$ cellula ibrida. Dopo circa 40 minuti le proteine si mescolano sull'intera superficie: dimostra che il doppio strato di fosfolipidi è \textbf{fluido} e le proteine si muovono al suo interno.

\rubrica{Movimenti nella membrana}
La disposizione ordinata dei fosfolipidi rende la membrana un \textit{cristallo liquido}; le catene idrocarburiche sono in costante movimento e ogni molecola può spostarsi lateralmente sulla stessa faccia del doppio strato. Movimenti dei fosfolipidi:
\begin{itemize}
  \item \textbf{flessione}: $\sim 10^{-9}$ s;
  \item \textbf{spostamento laterale} (nello stesso foglietto): $\sim 10^{-6}$ s, rapido;
  \item \textbf{flip-flop} (da un foglietto all'altro): $\sim 10^{5}$ s, molto lento.
\end{itemize}

% ---------------------------------------------------------------------------
\subsection{Tecnica del freeze-fracture}

Congelamento-frattura: le cellule sono congelate rapidamente in azoto liquido e poi fratturate con un colpo secco, così da separare le due metà del doppio strato lipidico e permettere l'analisi dell'interno della membrana al microscopio elettronico. Le particelle visibili sul lato interno esposto sono \textbf{proteine integrali}. Si distinguono la \textbf{faccia E} (esterna) e la \textbf{faccia P} (protoplasmatica).

% ---------------------------------------------------------------------------
\subsection{Modello a mosaico fluido}

Schema complessivo della membrana:
\begin{itemize}
  \item doppio strato fosfolipidico (teste idrofiliche verso l'esterno, code idrofobiche all'interno);
  \item colesterolo, inserito tra i fosfolipidi;
  \item glicolipidi e glicoproteine, con i gruppi glucidici rivolti verso l'esterno della cellula;
  \item proteine integrali (tra cui le proteine canale / di trasporto) e proteine periferiche;
  \item microfilamenti del citoscheletro collegati alla membrana.
\end{itemize}

\section{La membrana plasmatica: funzioni}

% ---------------------------------------------------------------------------
\subsection{Funzioni generali}

\begin{itemize}
  \item protettiva;
  \item forma;
  \item scambio di sostanze nutritive;
  \item trasporto;
  \item reattività agli stimoli esterni;
  \item risposta agli ormoni;
  \item interazione e comunicazione tra le cellule;
  \item proprietà antigeniche e reazioni immunitarie.
\end{itemize}

% ---------------------------------------------------------------------------
\subsection{Differenziazione funzionale (cellula epiteliale)}

\begin{itemize}
  \item \textbf{membrana apicale}: regolazione dell'ingresso di materiali nutritizi e di acqua; regolazione della secrezione; protezione;
  \item \textbf{membrana laterale}: contatto e adesione fra cellule; comunicazione cellulare;
  \item \textbf{membrana basale}: contatto con il substrato; creazione di gradienti ionici.
\end{itemize}

% ---------------------------------------------------------------------------
\subsection{Funzioni delle proteine di membrana}

\begin{itemize}
  \item \textbf{ancoraggio}: alcune proteine (es.\ le integrine) ancorano la cellula alla matrice extracellulare e si connettono ai microfilamenti intracellulari;
  \item \textbf{trasporto passivo}: certe proteine formano canali per il passaggio selettivo di ioni o molecole;
  \item \textbf{trasporto attivo}: alcune proteine pompano i soluti attraverso la membrana, con apporto diretto di energia;
  \item \textbf{attività enzimatica}: molti enzimi legati alla membrana catalizzano reazioni all'interno o sulla superficie;
  \item \textbf{trasduzione del segnale}: alcuni recettori legano molecole segnale (es.\ ormoni) e trasmettono l'informazione all'interno della cellula;
  \item \textbf{riconoscimento cellulare}: alcune glicoproteine fungono da marcatori di identificazione (es.\ gli antigeni di superficie delle cellule batteriche, riconosciuti come estranei);
  \item \textbf{giunzione intercellulare}: le proteine di adesione legano le membrane di cellule adiacenti.
\end{itemize}

% ---------------------------------------------------------------------------
\subsection{Diffusione e osmosi}

\rubrica{Diffusione}
Movimento delle molecole da una zona a maggiore concentrazione a una a minore concentrazione, fino a distribuzione uniforme (es.\ una zolletta di zucchero immersa in acqua si dissolve e diffonde uniformemente).

\rubrica{Osmosi}
Movimento netto di acqua attraverso una membrana selettivamente permeabile (permeabile all'acqua ma non al soluto). La \textbf{pressione osmotica} della soluzione è pari alla forza che occorre applicare per impedire l'ascesa del livello del fluido.
\begin{itemize}
  \item soluzione \textbf{ipotonica}: guadagno netto di acqua $\rightarrow$ la cellula si gonfia;
  \item soluzione \textbf{ipertonica}: perdita di acqua $\rightarrow$ la cellula si raggrinzisce;
  \item soluzione \textbf{isotonica}: né guadagno né perdita di acqua.
\end{itemize}

% ---------------------------------------------------------------------------
\subsection{Meccanismi di trasporto attraverso la membrana}

\begin{itemize}
  \item \textbf{trasporto passivo} (secondo gradiente, senza energia):
    \begin{itemize}
      \item non mediato: diffusione semplice (anche attraverso un canale acquoso);
      \item mediato: diffusione facilitata secondo gradiente di concentrazione;
    \end{itemize}
  \item \textbf{trasporto attivo}: contro gradiente, con apporto diretto di energia.
\end{itemize}

\rubrica{Trasporto passivo: diffusione semplice}
Piccole molecole (es.\ $O_2$, $CO_2$) attraversano direttamente il doppio strato lipidico.

\rubrica{Trasporto passivo: diffusione facilitata}
\begin{itemize}
  \item \textbf{proteine canale}: formano un poro selettivo (es.\ acquaporina, canale per l'acqua);
  \item \textbf{proteine carrier} (trasportatori): legano il soluto e cambiano conformazione esponendo il sito di legame alternativamente verso l'esterno o verso l'interno (legame $\rightarrow$ trasporto $\rightarrow$ dissociazione $\rightarrow$ ritorno alla conformazione iniziale); es.\ il trasportatore del glucosio GLUT1.
\end{itemize}

\rubrica{Trasporto attivo}
Contro gradiente, con consumo di energia (ATP).
\begin{itemize}
  \item \textbf{pompa $Na^+$--$K^+$}: ciclo conformazionale che, idrolizzando ATP, espelle $Na^+$ e importa $K^+$;
  \item \textbf{trasporto attivo secondario}: $Na^+$ e glucosio si legano insieme alla proteina carrier; la proteina cambia conformazione e li rilascia nel citosol. Il glucosio è trasportato contro il proprio gradiente sfruttando il gradiente del sodio.
\end{itemize}

% ---------------------------------------------------------------------------
\subsection{Trasporto di massa: endocitosi ed esocitosi}

\rubrica{Esocitosi}
Una vescicola secretoria si avvicina alla membrana plasmatica e si fonde con essa $\rightarrow$ rilascia il proprio contenuto nell'ambiente extracellulare; le proteine inserite nella membrana della vescicola diventano parte integrante della membrana plasmatica.

\rubrica{Fagocitosi}
Pieghe (pseudopodi) della membrana circondano la particella da ingerire formando un vacuolo (fagosoma) $\rightarrow$ il vacuolo, in seguito a una strozzatura, si libera all'interno della cellula $\rightarrow$ i lisosomi si fondono con esso (fagolisosoma) e riversano i loro enzimi digestivi sul materiale ingerito. Passaggi: intrappolamento $\rightarrow$ inglobamento $\rightarrow$ digestione $\rightarrow$ assorbimento. Nei mammiferi alcuni globuli bianchi (\textit{fagociti}) svolgono un processo simile.

\rubrica{Pinocitosi (endocitosi di massa)}
Goccioline di fluido sono intrappolate da pieghe della membrana plasmatica $\rightarrow$ subiscono una strozzatura e divengono piccole vescicole piene di fluido $\rightarrow$ il contenuto viene trasferito al citosol.

\rubrica{Endocitosi mediata da recettore}
Specifiche macromolecole si legano a recettori situati nelle fossette rivestite (di clatrina). Esempio: assorbimento delle LDL.
\begin{enumerate}
  \item la LDL si attacca a recettori specifici nelle fossette rivestite della membrana plasmatica;
  \item l'endocitosi forma una vescicola rivestita nel citosol; il rivestimento viene subito rimosso;
  \item la vescicola non rivestita si fonde con un endosoma;
  \item i recettori vengono riciclati e tornano sulla membrana plasmatica;
  \item la vescicola contenente le LDL si fonde con un lisosoma (lisosoma secondario); gli enzimi idrolitici rilasciano dalle LDL il colesterolo, utilizzato dalla cellula.
\end{enumerate}

\section{Il citoscheletro}

% ---------------------------------------------------------------------------
\subsection{Struttura e funzioni}

Il \textbf{citoscheletro} è un sistema di filamenti proteici; è una struttura \textbf{dinamica}, continuamente scomposta e riassemblata.

Funzioni:
\begin{itemize}
  \item mantiene la struttura e la forma della cellula (\textbf{struttura e supporto});
  \item fornisce ancoraggio per organuli e macromolecole (\textbf{organizzazione spaziale});
  \item permette il \textbf{trasporto intracellulare};
  \item nella divisione cellulare forma le fibre del fuso mitotico;
  \item permette il movimento cellulare (\textbf{contrattilità e motilità}).
\end{itemize}

Tre componenti (reti di fibre di diverso tipo):
\begin{itemize}
  \item \textbf{microtubuli}: $\phi$ 20--25 nm;
  \item \textbf{filamenti intermedi}: $\phi$ 8--12 nm;
  \item \textbf{microfilamenti}: $\phi$ 5--7 nm.
\end{itemize}

% ---------------------------------------------------------------------------
\subsection{Microtubuli}

\begin{itemize}
  \item diametro 20--25 nm; costituiti da subunità proteiche di \textbf{tubulina} ($\alpha$ e $\beta$);
  \item agiscono da impalcatura per determinare la forma della cellula;
  \item permettono il movimento degli organelli e delle vescicole;
  \item formano le fibre del fuso per separare i cromosomi durante la mitosi;
  \item all'interno di ciglia e flagelli permettono il movimento cellulare.
\end{itemize}

\rubrica{Formazione e polarità}
\begin{itemize}
  \item si formano per aggiunta di eterodimeri di $\alpha$-tubulina e $\beta$-tubulina a un'estremità del cilindro cavo;
  \item il cilindro possiede una \textbf{polarità}: a un'estremità (estremità $+$) si aggiungono dimeri, all'estremità opposta (estremità $-$) si rimuovono;
  \item per ogni giro di spirale sono necessari 13 dimeri.
\end{itemize}

\rubrica{Proteine associate ai microtubuli (MAP)}
\begin{itemize}
  \item presentano 3 siti di legame alla tubulina, intervallati da una distanza sufficiente a permettere il legame a 3 subunità separate;
  \item le code delle MAP si proiettano all'esterno, dove possono interagire con altri componenti cellulari.
\end{itemize}

\rubrica{Proteine motrici: la chinesina}
Proteina motrice dei microtubuli.
\begin{itemize}
  \item struttura ($\sim$80 nm): 2 catene pesanti avvolte a spirale nella regione a stelo e 2 catene leggere legate alla parte terminale delle catene pesanti; regioni: teste (sito catalitico), collo, stelo, coda;
  \item meccanismo: la molecola ``cammina'' lungo il microtubulo alternando l'attacco e il rilascio del ``piede''; le due teste effettuano movimenti identici ma alternati (passo di 16 nm).
\end{itemize}

\rubrica{Proteine motrici: la dineina citoplasmatica}
Contiene due catene pesanti e alcune catene leggere e intermedie; ciascuna catena pesante ha una grande testa globulare generante forza, un peduncolo con sito di legame per il microtubulo e uno stelo.
\begin{itemize}
  \item lungo lo stesso microtubulo le due proteine muovono le vescicole in direzioni opposte: la \textbf{chinesina} verso l'estremità $+$ (movimento anterogrado), la \textbf{dineina} verso l'estremità $-$ (movimento retrogrado);
  \item le proteine motrici sono unite alla membrana della vescicola tramite un intermediario: la chinesina tramite proteine di membrana (integrali o periferiche), la dineina tramite un complesso di proteine solubili detto \textbf{dinactina};
  \item mediano il trasporto di vescicole, gruppi vescicolari-tubulari (VTC) e organelli.
\end{itemize}

% ---------------------------------------------------------------------------
\subsection{Filamenti intermedi}

\begin{itemize}
  \item provvedono alla resistenza meccanica delle cellule soggette a stress fisici (cellule muscolari, epiteliali, neuroni);
  \item sono spesso interconnessi ad altri filamenti del citoscheletro da sottili ponti trasversali proteici (es.\ la plectina);
  \item bastoncini flessibili di $\sim$10 nm di diametro, costituiti da \textit{protofilamenti}, a loro volta formati da subunità proteiche avvolte a elica.
\end{itemize}

\rubrica{Assemblaggio}
\begin{enumerate}
  \item ciascun monomero è costituito da domini terminali globulari separati da una lunga regione ad $\alpha$-elica;
  \item coppie di monomeri si associano parallelamente, con le estremità allineate, a formare dimeri (omodimeri o eterodimeri);
  \item i dimeri si associano in maniera antiparallela e sfalsata a formare tetrameri;
  \item i tetrameri si associano lateralmente per formare un'unità di filamento ($\sim$60 nm);
  \item i filamenti intermedi molto lunghi sono formati dall'associazione di tali unità di lunghezza.
\end{enumerate}

\rubrica{Nella cellula epiteliale}
I filamenti intermedi si irradiano per tutta la cellula, ancorati sia alla superficie esterna del nucleo sia alla superficie interna della membrana plasmatica. Le connessioni con il nucleo avvengono tramite proteine che attraversano entrambe le membrane dell'involucro nucleare; quelle con la membrana plasmatica tramite siti specializzati come i desmosomi e gli emidesmosomi.

% ---------------------------------------------------------------------------
\subsection{Microfilamenti}

\begin{itemize}
  \item diametro 5--7 nm; costituiti da due catene di molecole di actina intrecciate tra loro (doppia elica);
  \item composti dalle subunità globulari della proteina \textbf{actina};
  \item in presenza di ATP i monomeri di actina polimerizzano formando un filamento flessibile a elica.
\end{itemize}

\rubrica{Proteine che legano l'actina}
\begin{enumerate}
  \item \textbf{di nucleazione}: intervengono nel primo passo di formazione del filamento, riunendo nel giusto orientamento 2 o 3 monomeri di actina;
  \item \textbf{che sequestrano i monomeri}: impediscono la polimerizzazione dei monomeri in filamenti; regolano l'equilibrio monomero--polimero;
  \item \textbf{di incappucciamento}: regolano la lunghezza dei filamenti legandosi a una delle due estremità;
  \item \textbf{che polimerizzano i monomeri}: promuovono la crescita dei filamenti;
  \item \textbf{che depolimerizzano i filamenti}: consentono la depolimerizzazione e un rapido turnover nei siti di dinamici cambiamenti del citoscheletro (locomozione cellulare, fagocitosi, citocinesi);
  \item \textbf{che formano legami crociati}: alterano l'organizzazione tridimensionale di una popolazione di filamenti (fascicolazione);
  \item \textbf{che tagliano i filamenti}: si legano a un filamento esistente e lo rompono in due;
  \item \textbf{che legano la membrana}: si legano alla membrana plasmatica permettendone i movimenti.
\end{enumerate}

\rubrica{Microfilamenti nel microvillo}
I microvilli si trovano sulla superficie apicale di epiteli con funzione di assorbimento (rivestimento interno dell'intestino, parete dei tubuli renali). I filamenti di actina sono mantenuti in disposizione ordinata da proteine che formano fasci, come la \textbf{villina} e la \textbf{fimbrina}; la \textbf{miosina I} è presente tra la membrana plasmatica del microvillo e i filamenti periferici di actina.

% ---------------------------------------------------------------------------
\subsection{Confronto tra i tre componenti}

\begin{table}[htbp]
  \centering \footnotesize \renewcommand{\arraystretch}{1.25}
  \begin{tabularx}{\textwidth}{|l|>{\raggedright\arraybackslash}X|>{\raggedright\arraybackslash}X|>{\raggedright\arraybackslash}X|}
    \hline
     & \textbf{Microtubuli} & \textbf{Filamenti intermedi} & \textbf{Filamenti di actina} \\
    \hline
    Subunità incorporata & eterodimero di GTP--$\alpha\beta$-tubulina & circa 70 proteine diverse & monomeri di ATP-actina \\ \hline
    Sito di incorporazione & estremità $+$ ($\beta$-tubulina) & interno & estremità $+$ (sfrangiata) \\ \hline
    Polarità & sì & no & sì \\ \hline
    Attività enzimatica & GTPasi & nessuna & ATPasi \\ \hline
    Proteine motrici & chinesine, dineine & nessuna & miosine \\ \hline
    Proteine associate & MAP & plachine & proteine leganti actina \\ \hline
    Struttura & tubo rigido, cavo, inestensibile & filamenti resistenti, flessibili, estensibili & filamenti a elica, flessibili, inestensibili \\ \hline
    Dimensioni & diametro esterno 25 nm & diametro 10--12 nm & diametro 8 nm \\ \hline
    Distribuzione & tutti gli eucarioti & animali & tutti gli eucarioti \\ \hline
    Funzioni primarie & supporto, trasporto intracellulare, organizzazione della cellula & supporto strutturale & motilità, contrattilità \\ \hline
    Distribuzione sub-cellulare & citoplasma & citoplasma e nucleo & citoplasma \\ \hline
  \end{tabularx}
\end{table}

% ---------------------------------------------------------------------------
\subsection{Centrioli}

\begin{itemize}
  \item diametro $\approx$ 0,2\,$\mu$m; lunghezza $\approx$ 0,4\,$\mu$m;
  \item strutture cilindriche che contengono 9 fibrille regolarmente spaziate, ognuna formata da 3 microtubuli (arrangiamento 9$\times$3), localizzate vicino al nucleo;
  \item sono presenti quasi sempre in coppia, ad angolo retto l'uno rispetto all'altro;
  \item durante la divisione cellulare si duplicano, dando origine a 2 coppie di centrioli che si spostano ai poli opposti della cellula, dove funzionano come centri per l'organizzazione dei microtubuli che formano il fuso mitotico.
\end{itemize}

% ---------------------------------------------------------------------------
\subsection{Centrosoma}

Struttura complessa che contiene 2 centrioli circondati da materiale pericentriolare (PCM), amorfo e opaco agli elettroni, in cui avviene la nucleazione dei microtubuli.
\begin{itemize}
  \item è il principale sito di inizio dei microtubuli (MTOC) e si trova tipicamente al centro della cellula, in una fitta rete di microtubuli;
  \item si duplica una volta per ciclo cellulare, in coordinazione con la replicazione del DNA: ciascun centriolo genera un centriolo figlio, e alla mitosi le due coppie migrano ai poli opposti.
\end{itemize}

% ---------------------------------------------------------------------------
\subsection{Ciglia e flagelli}

\begin{itemize}
  \item sottili organelli motori, simili a peli, che si proiettano dalla superficie di vari tipi di cellule;
  \item sono sottili filamenti avvolti da un'estroflessione della membrana plasmatica;
  \item alla base sono ancorati dal centriolo basale (corpo basale), da cui si irradia in profondità la radice del ciglio;
  \item le \textbf{ciglia} sono più corte e più numerose rispetto ai \textbf{flagelli}.
\end{itemize}

\rubrica{Battito}
\begin{itemize}
  \item le ciglia alternano una \textit{fase attiva} (di spinta) e una \textit{fase di ritorno}; il battito, flessibile, inizia alla base e si propaga verso l'apice, con ritmo \textit{metacronale} (le ciglia di file adiacenti sono in stadi diversi del battito);
  \item meccanismo: i bracci di dineina muovono i microtubuli formando e rompendo i legami con i microtubuli adiacenti, così che ciascuno ``cammina'' lungo quello vicino; le proteine di connessione (nexina) impediscono lo scivolamento oltre un certo limite $\rightarrow$ l'azione motrice si traduce nella flessione dei microtubuli (movimento oscillatorio);
  \item i \textbf{flagelli} si muovono con una forma d'onda (asimmetrica nel caso dell'alga \textit{Chlamydomonas}); es.\ il flagello degli spermatozoi.
\end{itemize}

% ---------------------------------------------------------------------------
\subsection{Assonema}

Struttura interna di ciglia e flagelli: microtubuli raggruppati in 9 coppie (doppiette) disposte alla periferia più una coppia al centro $\rightarrow$ complesso ``9$+$2''.

\rubrica{Sezione trasversale}
\begin{itemize}
  \item 9 doppiette esterne di microtubuli, ciascuna formata da un \textbf{tubulo A} e un \textbf{tubulo B}, più una coppia centrale racchiusa nella guaina centrale;
  \item \textbf{bracci di dineina} (interno ed esterno) sul tubulo A;
  \item \textbf{raggi} (spoke) diretti verso il centro;
  \item \textbf{ponte tra le coppie} (nexina), che collega le doppiette adiacenti.
\end{itemize}

\rubrica{Sezione longitudinale}
Raggi, coppie periferiche e guaina centrale; i bracci di dineina si ripetono lungo l'asse con periodo di 96 nm (spaziature 24, 32, 40, 24 nm).

\rubrica{Corpo basale}
Alla base dell'assonema; deriva dai centrioli (arrangiamento 9$\times$3 di triplette) e ancora ciglia e flagelli alla cellula.

\section{Le giunzioni cellulari}

% ---------------------------------------------------------------------------
\subsection{La matrice extracellulare}

\textbf{Matrice extracellulare} (ECM): rete organizzata di materiale extracellulare presente nelle immediate vicinanze della membrana plasmatica.

\rubrica{Composizione}
\begin{itemize}
  \item \textbf{collagene}: proteina fibrosa;
  \item \textbf{fibronectine}: glicoproteine della ECM che si legano alle integrine e ad altri recettori presenti nella membrana plasmatica;
  \item \textbf{laminina};
  \item \textbf{proteoglicano}: enormi complessi proteina--polisaccaride che occupano la maggior parte del volume dello spazio extracellulare, con siti di legame reciproco e per recettori (integrine) localizzati sulla superficie delle cellule;
  \item \textbf{integrina}: recettore di membrana che lega la ECM al citoscheletro (filamento intermedio, microfilamenti).
\end{itemize}

\rubrica{Caratteristiche}
\begin{itemize}
  \item le proteine presenti nella matrice extracellulare servono come impalcature, travi, cemento e reti;
  \item può assumere diverse forme nei differenti tessuti;
  \item la maggior parte delle proteine presenti sono di tipo fibroso;
  \item può avere un ruolo regolatorio nel determinare la forma e le attività di una cellula.
\end{itemize}

% ---------------------------------------------------------------------------
\subsection{Lamina basale}

Funzioni:
\begin{itemize}
  \item fornisce un supporto meccanico per l'adesione cellulare;
  \item genera segnali per la sopravvivenza cellulare;
  \item serve da substrato per la migrazione cellulare;
  \item separa tessuti adiacenti all'interno di un organo;
  \item funziona da barriera al passaggio di macromolecole.
\end{itemize}

% ---------------------------------------------------------------------------
\subsection{Panoramica delle giunzioni cellulari}

Quattro tipi principali: \textbf{giunzione gap} (comunicante), \textbf{giunzioni aderenti}, \textbf{giunzioni strette}, \textbf{desmosoma}.
\begin{itemize}
  \item \textbf{giunzione gap}: connessoni;
  \item \textbf{giunzioni aderenti}: caderina, filamenti di actina;
  \item \textbf{giunzioni strette}: complessi proteici;
  \item \textbf{desmosoma}: cadherin, filamenti di actina.
\end{itemize}

% ---------------------------------------------------------------------------
\subsection{Desmosomi}

Ciascuna delle due cellule compartecipa alla formazione del desmosoma. Ciascun desmosoma è costituito da un paio di dischi a forma di bottone associati con le membrane plasmatiche di cellule adiacenti e da filamenti proteici intercellulari che connettono queste cellule. I filamenti intermedi attaccati a questi dischi sono connessi con altri desmosomi.

Poiché sono ancorati su entrambi i lati cellulari, risultano essere molto robusti.

\rubrica{Struttura}
\begin{itemize}
  \item struttura densa (\textbf{placca}) sul versante citoplasmatico della membrana;
  \item molecole di adesione (\textbf{desmogleina}, \textbf{desmocollina}) si estendono attraverso lo spazio intercellulare fino all'altra cellula, in corrispondenza delle proteine della placca;
  \item a livello citoplasmatico sono presenti proteine del citoscheletro, in particolare la cheratina;
  \item nella placca: \textbf{placoglobina} e \textbf{desmoplachina}, a cui si ancora il filamento intermedio.
\end{itemize}

% ---------------------------------------------------------------------------
\subsection{Emidesmosomi}

Siti differenziati alla superficie basale delle cellule epiteliali, dove le cellule sono attaccate alla membrana basale sottostante.

\begin{itemize}
  \item le molecole di \textbf{integrina} $\alpha_6\beta_4$ delle cellule epidermiche sono collegate ai filamenti intermedi citoplasmatici da una proteina chiamata \textbf{plectina}, presente nelle placche dense, e alla membrana basale da filamenti di ancoraggio formati da un tipo particolare di laminina (\textit{laminina-5});
  \item negli emidesmosomi è presente anche una seconda proteina transmembrana (\textbf{BP180});
  \item le fibre di collagene (collagene tipo VII) fanno parte del derma sottostante.
\end{itemize}

% ---------------------------------------------------------------------------
\subsection{Le giunzioni strette (tight junctions)}

\begin{itemize}
  \item le membrane plasmatiche adiacenti giungono in intimo contatto tra loro, annullando lo spazio intercellulare;
  \item lungo la zona di contatto gli strati esterni delle 2 membrane si fondono, formando un'unica struttura \textbf{pentalaminare} data dalla fusione di proteine intrinseche di membrana;
  \item una giunzione serrata si forma per mezzo di connessioni tra file di proteine di membrana adiacenti; tali proteine sono strettamente impacchettate in file che sigillano lo spazio intercellulare, impedendo il passaggio di materiali negli spazi tra le cellule.
\end{itemize}

Impediscono il passaggio di materiali attraverso gli spazi intercellulari.

% ---------------------------------------------------------------------------
\subsection{Le giunzioni comunicanti (gap junctions)}

Permettono il trasferimento di piccole molecole e ioni tra cellule adiacenti.

\begin{itemize}
  \item connessioni di tipo proteico (\textbf{connessoni}) legano le due membrane plasmatiche;
  \item le particelle proteiche sono attraversate da un canale idrofilo che permette lo scambio di ioni e piccole molecole tra le due cellule;
  \item le due membrane plasmatiche contengono cilindri costituiti da sei molecole di \textbf{connessina}; i due cilindri di membrane opposte sono uniti a formare un canale che connette i compartimenti citoplasmatici delle due cellule;
  \item il connessone (cilindro) è formato da subunità di connessina;
  \item il poro di una giunzione comunicante può aprirsi e chiudersi (conformazioni \textit{chiuso}/\textit{aperto});
  \item le proteine sono distribuite a distanze regolari, formando tra le parti sporgenti una rete di spazi che separa le due membrane adiacenti ed è accessibile dagli interstizi intercellulari circostanti.
\end{itemize}

% ---------------------------------------------------------------------------
\subsection{Le giunzioni aderenti}

L'actina e le caderine permettono alle membrane di essere strettamente legate in modo da impedirne la disgiunzione.

\begin{itemize}
  \item \textbf{caderina}: proteina transmembrana di adesione, extracellulare;
  \item \textbf{$\alpha$-actinina}: lega il filamento di actina alla $\beta$-catenina;
  \item la caderina si lega, sul versante citoplasmatico, a \textbf{$\beta$-catenina}, \textbf{$\alpha$-catenina} e alla \textbf{catenina p120};
  \item la $\beta$-catenina è collegata a una proteina che lega l'actina, ancorando il complesso al filamento di actina;
  \item i segnali associati al complesso di giunzione sono trasmessi al citoplasma e al nucleo.
\end{itemize}

% ---------------------------------------------------------------------------
\subsection{I plasmodesmi}

Mettono in comunicazione il citoplasma di cellule adiacenti, permettendo il passaggio di acqua, ioni e piccole molecole.

\begin{itemize}
  \item internamente i canali sono rivestiti dalle membrane plasmatiche fuse delle due cellule adiacenti;
  \item il \textbf{desmotubulo} consiste in una membrana in continuità con il reticolo endoplasmatico (RE) di entrambi i lati della membrana plasmatica;
  \item l'\textbf{annulus} è la regione anulare del plasmodesma attorno al desmotubulo;
  \item passano attraverso l'\textit{annulus} le molecole che si spostano da una cellula all'altra;
  \item il trasporto di proteine e mRNA attraverso i plasmodesmi è coinvolto nella localizzazione tissutale di alcuni fattori di trascrizione.
\end{itemize}

\section{Le endomembrane}

Il sistema delle \textbf{endomembrane} comprende: reticolo endoplasmatico, complesso di Golgi, lisosomi, perossisomi, vacuoli.

% ---------------------------------------------------------------------------
\subsection{Reticolo endoplasmatico: struttura}

\rubrica{Liscio (REL)}
\begin{itemize}
  \item costituito da membrane tubulari;
  \item sistema di canali interconnessi che si diramano nel citoplasma.
\end{itemize}

\rubrica{Rugoso (RER)}
\begin{itemize}
  \item costituito da sacche appiattite (\textit{cisterne});
  \item presenza di ribosomi sulla superficie esterna;
  \item è continuo con la membrana esterna dell'involucro nucleare.
\end{itemize}

% ---------------------------------------------------------------------------
\subsection{Reticolo endoplasmatico liscio: funzioni}

\begin{itemize}
  \item biogenesi di membrane;
  \item sintesi di colesterolo e ormoni steroidei;
  \item nel fegato: detossificazione di composti organici;
  \item sequestro di ioni calcio nelle cellule muscolari.
\end{itemize}

% ---------------------------------------------------------------------------
\subsection{Reticolo endoplasmatico rugoso: funzioni}

\begin{itemize}
  \item \textbf{sintesi delle proteine}: secrete dalla cellula, integrali di membrana, solubili situate all'interno di organuli cellulari;
  \item sintesi delle catene di carboidrati;
  \item sintesi dei fosfolipidi.
\end{itemize}

\rubrica{Sintesi di una proteina di secrezione o di un enzima lisosomale}
\begin{enumerate}
  \item la sintesi inizia su un ribosoma libero; quando la sequenza segnale emerge dal ribosoma si lega alla \textbf{SRP} (particella di riconoscimento del segnale), che blocca la traduzione;
  \item il complesso catena nascente--ribosoma--SRP prende contatto con la membrana del RE mediante un recettore per SRP;
  \item rilascio dell'SRP e associazione del ribosoma con il \textbf{traslocone};
  \item il peptide segnale si lega all'interno del traslocone, spostando il tappo del canale e permettendo al polipeptide di attraversare la membrana; quando il polipeptide nascente è entrato nel lume del RE, il peptide segnale viene tagliato da una proteina di membrana e la proteina si ripiega in modo corretto con l'aiuto di proteine specifiche.
\end{enumerate}

\rubrica{Sintesi di una proteina integrale di membrana}
Modello per una proteina con un singolo segmento transmembrana vicino all'estremità N-terminale del polipeptide nascente.
\begin{itemize}
  \item il polipeptide nascente entra nel traslocone come se fosse una proteina di secrezione; l'entrata della sequenza idrofobica transmembrana blocca l'ulteriore traslocazione del polipeptide nascente attraverso il canale;
  \item il traslocone, aprendosi lateralmente, espelle il segmento transmembrana nel doppio strato lipidico;
  \item nella riorientazione il traslocone dispone il segmento transmembrana attraverso le interazioni con le cariche opposte (positive e negative) presenti sulla sua superficie;
  \item il traslocone si apre lateralmente ed espelle il segmento transmembrana nel bilayer, con la disposizione finale della proteina: estremità C-terminale nel lume, estremità N-terminale nel citosol (o viceversa, a seconda dell'orientamento).
\end{itemize}

\rubrica{Mantenimento dell'asimmetria della membrana}
Quando una proteina è sintetizzata nel RER, viene subito inserita nel doppio strato lipidico in un orientamento prevedibile determinato dalla sua sequenza amminoacidica. Questo orientamento è mantenuto durante il suo viaggio nel sistema delle endomembrane. Le catene di carboidrati, aggiunte per prime nel RE, forniscono un modo conveniente per determinare l'asimmetria della membrana, perché sono sempre presenti sul lato cisternale delle membrane citoplasmatiche, che diventa il lato esoplasmico della membrana plasmatica dopo la fusione delle vescicole con la membrana cellulare.

% ---------------------------------------------------------------------------
\subsection{Complesso di Golgi: struttura}

\begin{itemize}
  \item localizzato tra il reticolo endoplasmatico e la membrana plasmatica;
  \item cisterne appiattite discoidali, curve, disposte in pile ordinate e interconnesse da tubuli membranosi;
  \item bordi dilatati associati a vescicole che gemmano dalla parte periferica di ogni cisterna.
\end{itemize}

Organizzazione delle cisterne, da un polo all'altro:
\[
\text{rete } cis \text{ del Golgi (CGN)} \rightarrow \text{cisterne } cis \rightarrow \text{cisterne mediali} \rightarrow \text{cisterne } trans \rightarrow \text{rete } trans \text{ del Golgi (TGN)}.
\]

% ---------------------------------------------------------------------------
\subsection{Complesso di Golgi: funzioni}

\begin{itemize}
  \item \textbf{modificazioni chimiche delle proteine}: rimozione di porzioni di sequenza amminoacidica; aggiunta di gruppi funzionali; aggiunta di unità lipidiche o glucidiche;
  \item \textbf{regolazione del movimento delle proteine}: proteine di secrezione, proteine di membrana, proteine destinate ad altri organuli (es.\ lisosomi).
\end{itemize}

% ---------------------------------------------------------------------------
\subsection{Trasporto delle proteine all'interno della cellula}

\begin{enumerate}
  \item i neopolipeptidi entrano nel RER;
  \item aggiunta di carboidrati a formare glicoproteine;
  \item vescicole di trasporto le trasferiscono sulla superficie \textit{cis} del complesso di Golgi;
  \item ulteriori modificazioni;
  \item le glicoproteine si portano sulla superficie \textit{trans}, dove sono impacchettate su vescicole di trasporto;
  \item vanno sulla membrana plasmatica (o su altri organuli);
  \item il contenuto della vescicola è rilasciato dalla cellula.
\end{enumerate}

% ---------------------------------------------------------------------------
\subsection{Lisosomi}

\begin{itemize}
  \item organelli digestivi;
  \item contengono $\sim$50 enzimi idrolitici: fosfatasi, nucleasi, proteasi, lipasi acida e fosfolipasi, polisaccaridasi e oligosaccaridasi;
  \item attività ottimale a pH acido;
  \item pompa protonica sulla membrana.
\end{itemize}

\rubrica{Funzioni}
\begin{itemize}
  \item \textbf{degradazione delle sostanze} che arrivano dall'ambiente esterno;
  \item \textbf{turnover degli organelli}:
    \begin{enumerate}
      \item un organello viene circondato da una membrana fornita da una cisterna del RE;
      \item si fonde con un lisosoma a formare un \textit{autofagolisosoma};
      \item l'organello viene degradato e i prodotti resi disponibili per la cellula;
      \item al termine del processo l'organello diventa un \textit{corpo residuo};
      \item il corpo residuo può seguire due vie alternative: andare incontro a esocitosi, oppure rimanere nella cellula indefinitamente sotto forma di \textit{granulo di lipofuscina}.
    \end{enumerate}
\end{itemize}

\rubrica{Autofagia}
Un mitocondrio e un perossisoma vengono racchiusi in un involucro a doppia membrana derivato dal RE: si forma così un vacuolo autofagico (autofagolisosoma), che deve fondersi con un lisosoma perché il suo contenuto sia digerito.

% ---------------------------------------------------------------------------
\subsection{Perossisomi}

Detti anche \textit{microcorpi}.
\begin{itemize}
  \item vescicole circondate da membrana;
  \item $\phi$ 0,1--1,0\,$\mu$m;
  \item contengono enzimi ossidativi;
  \item sito di formazione e degradazione di $H_2O_2$;
  \item si duplicano per scissione da organelli preesistenti.
\end{itemize}

\rubrica{Meccanismo redox}
I perossisomi contengono enzimi che riducono l'ossigeno molecolare ad acqua in due tappe:
\begin{itemize}
  \item prima tappa (\textbf{ossidasi}): rimuove elettroni da substrati differenti ($RH_2$, es.\ acido urico o amminoacidi) $\rightarrow$ $R_2 + H_2O_2$;
  \item seconda tappa (\textbf{catalasi}): converte in acqua il perossido di idrogeno formatosi nella prima tappa $\rightarrow$ $2H_2O_2 \rightarrow O_2 + 2H_2O$.
\end{itemize}

% ---------------------------------------------------------------------------
\subsection{Vacuolo delle cellule vegetali}

\begin{itemize}
  \item peculiare delle cellule vegetali;
  \item occupa il 90\% della cellula;
  \item delimitato da una singola membrana: il \textit{tonoplasto};
  \item all'interno alta concentrazione di ioni $\rightarrow$ l'acqua entra per osmosi $\rightarrow$ turgore;
  \item siti di digestione intracellulare;
  \item pH acido per la presenza di una pompa protonica sul tonoplasto.
\end{itemize}

\section{Mitocondri e cloroplasti}

I mitocondri e i cloroplasti sono quegli organuli cellulari in grado di convertire l'energia in forme utilizzabili dalle cellule.

\rubrica{Ciclo dell'energia}
\[
\text{glicolisi} \rightarrow \text{fosforilazione ossidativa} \rightarrow ATP \rightarrow \left\{\text{sintesi proteine, energia meccanica, trasporto attivo, calore}\right\} \rightarrow ADP \rightarrow
\]

\rubrica{Respirazione aerobica e fotosintesi a confronto}
\begin{itemize}
  \item \textbf{respirazione aerobica} (mitocondri, nella maggior parte delle cellule eucariotiche): glucosio $+ O_2 \rightarrow ATP + CO_2 + H_2O$;
  \item \textbf{fotosintesi} (cloroplasti, in alcune cellule vegetali e algali): luce $+ CO_2 + H_2O \rightarrow ATP + O_2 + $ glucosio.
\end{itemize}
I prodotti dell'una sono i substrati dell'altra: il ciclo si chiude tra le due vie.

% ---------------------------------------------------------------------------
\subsection{Mitocondri}

Sono gli organuli respiratori della cellula, deputati alla produzione dell'energia necessaria per molte funzioni cellulari.

\rubrica{Caratteristiche generali}
\begin{itemize}
  \item forma di bastoncino; $\phi$ 0,2\,$\mu$m, lunghezza 2--4\,$\mu$m;
  \item composizione lipoproteica: 65--70\% proteine, 25--30\% lipidi;
  \item presenza di piccole quantità di RNA e DNA;
  \item contengono un elevato contenuto di enzimi capaci di ossidare i prodotti dell'assorbimento intestinale, degradandoli ad $H_2O$ e $CO_2$, liberando energia che viene convertita in energia chimica mediante la fosforilazione di adenosin-di-fosfato (ADP) in adenosin-tri-fosfato (ATP): $ADP + P \rightarrow ATP$.
\end{itemize}

% ---------------------------------------------------------------------------
\subsection{Struttura}

Forma allungata; circoscritti da due membrane:
\begin{itemize}
  \item \textbf{membrana mitocondriale esterna}: decorso regolare, racchiude i mitocondri;
  \item \textbf{membrana mitocondriale interna}: si solleva in numerose pliche denominate \textbf{creste mitocondriali}, orientate trasversalmente all'asse maggiore del mitocondrio.
\end{itemize}
Spazio tra le due membrane $\rightarrow$ \textbf{spazio intermembrana}. Spazio all'interno del mitocondrio $\rightarrow$ \textbf{matrice mitocondriale}.

Sulla membrana interna: particelle di \textbf{ATP sintasi}. Nella matrice: DNA mitocondriale, ribosomi, granuli.

% ---------------------------------------------------------------------------
\subsection{La respirazione mitocondriale}

Le diverse fasi della respirazione aerobica hanno luogo in sedi specifiche:
\begin{itemize}
  \item \textbf{citoplasma}: glicolisi (glucosio $\rightarrow$ piruvato) $\rightarrow$ 2 ATP;
  \item \textbf{ingresso nel mitocondrio}: formazione dell'acetil coenzima A (piruvato $\rightarrow$ acetil-CoA);
  \item \textbf{matrice mitocondriale}: ciclo dell'acido citrico (ciclo di Krebs) $\rightarrow$ 2 ATP;
  \item \textbf{creste mitocondriali}: trasporto degli elettroni e chemiosmosi (catena respiratoria) $\rightarrow$ 32 ATP.
\end{itemize}

% ---------------------------------------------------------------------------
\subsection{Biochimica}

Tre classi di enzimi agiscono in sequenza.

\rubrica{Nella matrice mitocondriale}
\textbf{Enzimi ossidativi del ciclo di Krebs}: degradano i prodotti derivanti dall'assorbimento intestinale ad anidride carbonica ($CO_2$), liberando elettroni (o i loro equivalenti atomi di idrogeno, $H^+$).

\rubrica{Sulla membrana mitocondriale interna}
\begin{itemize}
  \item \textbf{enzimi della catena respiratoria} (sistema di trasporto degli elettroni): trasportano coppie di elettroni (o i loro equivalenti ioni $H^+$), liberati durante il ciclo di Krebs, attraverso varie reazioni intermedie al termine delle quali gli elettroni si combinano con ossigeno molecolare formando acqua ($H_2O$) e cedendo energia;
  \item \textbf{enzimi fosforilativi} (ATP sintasi): catalizzano la sintesi di adenosintrifosfato (ATP) a partire da adenosindifosfato (ADP) e dal fosfato, utilizzando l'energia ceduta da una coppia di elettroni.
\end{itemize}

La catena respiratoria è organizzata in complessi (I, II, III, IV) sulla membrana interna; il gradiente protonico generato attraversa il complesso V (ATP sintasi), producendo ATP a partire da ADP e $P_i$.

% ---------------------------------------------------------------------------
\subsection{Altre funzioni}

\begin{itemize}
  \item metabolismo dei lipidi e dei fosfolipidi;
  \item capacità di ossidare gli acidi grassi;
  \item partecipazione, insieme al reticolo endoplasmatico liscio, alla biosintesi degli ormoni steroidei;
  \item in grado di accumulare e di concentrare varie sostanze;
  \item nella matrice: modesta sintesi proteica e presenza di piccole quantità di DNA, di ribosomi e delle varie classi di RNA.
\end{itemize}

% ---------------------------------------------------------------------------
\subsection{Genoma mitocondriale}

\begin{itemize}
  \item presenza di DNA e di tutto l'apparato per la sintesi proteica: ribosomi, tRNA, mRNA;
  \item DNA mitocondriale simile a quello batterico;
  \item si formano per divisione di mitocondri preesistenti;
  \item si duplicano come i batteri;
  \item le molecole di DNA mitocondriale sono troppo piccole per codificare tutte le proteine di questo organulo $\rightarrow$ i mitocondri sono quindi \textbf{semi-autonomi}, cioè capaci di duplicarsi e presiedere alla sintesi di alcuni costituenti chimici.
\end{itemize}

\rubrica{Analogie con i batteri}
\begin{itemize}
  \item il DNA ha forma circolare e sono assenti le proteine cromosomiche;
  \item i ribosomi sono di grandezza simile o inferiore a quelli dei batteri;
  \item la membrana interna del mitocondrio con il sistema di trasporto degli elettroni è paragonabile alla membrana plasmatica dei batteri, che ha funzione respiratoria e che forma invaginazioni (\textit{mesosomi}) simili alle creste mitocondriali;
  \item sulla base di queste analogie, si ipotizza che i mitocondri possano rappresentare batteri assunti nel corso dell'evoluzione all'interno di cellule primitive e che abbiano realizzato una forma di simbiosi con la cellula ospite.
\end{itemize}

% ---------------------------------------------------------------------------
\subsection{Cloroplasti}

\rubrica{Struttura}
Due membrane:
\begin{itemize}
  \item \textbf{membrana esterna}: porine, $\phi$ 1 nm;
  \item \textbf{membrana interna}: impermeabile (presenza di trasportatori), ripiegata a formare i \textbf{tilacoidi}.
\end{itemize}
\begin{itemize}
  \item i \textbf{tilacoidi} sono disposti in pile ordinate dette \textbf{grana};
  \item spazio interno ai tilacoidi $\rightarrow$ \textbf{lume};
  \item spazio interno del cloroplasto (esterno ai tilacoidi) $\rightarrow$ \textbf{stroma}, che contiene gli enzimi per la sintesi dei carboidrati.
\end{itemize}

Presenza di DNA, residuo di un batterio endosimbionte ancestrale.

% ---------------------------------------------------------------------------
\subsection{Fotosintesi clorofilliana}

\rubrica{Fase luminosa}
Assorbimento dell'energia luminosa, che viene conservata sotto forma di ATP e NADPH. Si ha sviluppo di ossigeno.

\rubrica{Fase oscura}
ATP e NADPH vengono utilizzati per ridurre la $CO_2$ a formare glucosio e altri prodotti organici.

Reazione complessiva:
\[
6CO_2 + 6H_2O \rightarrow C_6H_{12}O_6 + 6O_2.
\]

\rubrica{Fotosistemi}
\begin{itemize}
  \item \textbf{fotosistema I}: complesso proteico (dimero P700) associato all'antenna LHCI; assorbe energia luminosa e, tramite una catena di trasportatori (tra cui la ferredossina), riduce il NADP\textsuperscript{+} a NADPH;
  \item \textbf{fotosistema II}: complesso proteico (dimero P680, proteine D1/D2) associato all'antenna LHCII; il \textbf{complesso sviluppante ossigeno} (con centro Mn--Ca) scinde l'acqua liberando elettroni, protoni e $O_2$;
  \item il trasporto di elettroni tra PSII e PSI avviene attraverso trasportatori mobili (plastochinone PQ, citocromo $b_6f$, plastocianina PC), con generazione di un gradiente protonico attraverso la membrana tilacoidale.
\end{itemize}

% ---------------------------------------------------------------------------
\subsection{Origine di mitocondri e cloroplasti}

Teoria dell'\textbf{endosimbiosi}: cellula ospite di tipo procariotico $\rightarrow$ acquisizione di batteri aerobi (endosimbionti), tramite numerose invaginazioni della membrana plasmatica.
\begin{itemize}
  \item i batteri aerobi diventano \textbf{mitocondri};
  \item il reticolo endoplasmatico e l'involucro nucleare si formano dall'invaginazione della membrana plasmatica (non incluso nell'ipotesi dell'endosimbiosi seriale);
  \item batteri fotosintetici diventano \textbf{cloroplasti}, in cellule eucariotiche vegetali e in alcuni protisti;
  \item le cellule eucariotiche di animali, funghi e alcuni protisti possiedono mitocondri ma non cloroplasti.
\end{itemize}


% =============================================================================
% CAPITOLO 3 — Ciclo cellulare e genetica
% =============================================================================
\chapter{Ciclo cellulare e genetica}
% \input{03_ciclo_cellulare_e_genetica/01_...}

% =============================================================================
% CAPITOLO 4 — Istologia
% =============================================================================
\chapter{Istologia}
% \input{04_istologia/01_...}

\end{document}
